\documentclass[12pt]{article}

\usepackage{amsmath,amssymb,hyperref,longtable}
\topmargin -.5cm
\textheight 21cm
\oddsidemargin -.125cm
\textwidth 17cm

\def\be{\begin{equation}}
\def\ee{\end{equation}}
\def\ie{\begin{equation}\begin{aligned}}
\def\fe{\end{aligned}\end{equation}}
\newcommand{\adss}{\ensuremath{AdS_5\times S^5}}
\newcommand{\D}{\ensuremath{\Delta}}
\newcommand{\Tr}{\mathrm{Tr}}

\begin{document}

\begin{center}
{\Large\bf Notes on Convenient Konishi Representatives for the Linearized SFT Equations}
\end{center}

\tableofcontents

\section{Goal}

The goal of this note is to identify useful representatives of the Konishi multiplet for the
linearized closed superstring field theory analysis around the \adss\ background.
The practical tension is the following.

\begin{itemize}
\item Table 1 of \texttt{references/1102.1209/1102.1209.tex} organizes the Konishi multiplet in terms of
the full $SO(2,4)\times SO(6)$ quantum numbers.
\item In the SFT construction of \texttt{references/2507.12921/2507.12921.tex}, the background is written in
flat ten-dimensional variables $X^\mu$ with only the split into $AdS_5$ and $S^5$ directions manifest, i.e.
effectively an $SO(1,4)\times SO(5)$-covariant description.
\end{itemize}

Therefore the most convenient state for the linearized equations is not automatically the most
useful state for seeing the expected Konishi anomalous dimension. In particular, odd orders in the
background expansion are RR and even orders are NSNS, and this often removes the obvious
constant correction for simple pure-sector trial states.

Throughout the note I will use the same background notation as in \texttt{draftKonishi.tex}:
\ie
\psi = F_{\alpha\beta}(X)V^{\alpha\beta}_{RR}
+ h_{\mu\nu}(X)V^{\mu\nu}_{NSNS}
+ \phi(X)D
+ A_\mu(X)V^\mu_{\rm aux.},
\fe
with
\ie
F^{(1)}_{\alpha\beta}\neq 0,\qquad
(h^{(2)}_{\mu\nu},\phi^{(2)},A^{(2)}_\mu)\neq 0,
\fe
so that the perturbative background alternates as
\be
\psi=\psi^{(1)}_{RR}+\psi^{(2)}_{NSNS}+\psi^{(3)}_{RR}+\cdots.
\ee

The massless effective linearized equation is
\be
Q\varphi+\sum_{n=1}^{\infty}\frac{1}{n!}[\psi^{\otimes n}\otimes\varphi]'=0.
\ee

\section{Konishi Data and Expected Strong-Coupling Behavior}

For a genuine member of the Konishi multiplet on the first excited string level, the strong-coupling
expansion reviewed in \texttt{1102.1209} takes the form
\be
\D = 2\lambda^{1/4}+b_0+\frac{b_1}{\lambda^{1/4}}+O(\lambda^{-1/2}),
\ee
with
\be
b_0=\D_0-4,
\ee
and, for the states that should belong to the same Konishi long multiplet,
\be
b_1=2.
\ee

In the radius conventions used in the SFT draft, one sets $\lambda^{1/4}=R$ (up to possible
scheme-dependent redefinitions of $R$). Then the expected Konishi behavior is
\be
\D = 2R+(\D_0-4)+\frac{2}{R}+O(R^{-2}).
\ee

The order counting relevant for the SFT computation is then:
\begin{itemize}
\item the piece $2R$ is the zeroth-order contribution;
\item the constant piece $\D_0-4$ is the first correction;
\item the $2/R$ term is the next correction.
\end{itemize}

The relevant expectations are then:
\begin{itemize}
\item bottom component, $\D_0=2$:
\be
\D_{\rm bottom}=2R-2+\frac{2}{R}+O(R^{-2});
\ee
\item level-4 descendants, $\D_0=4$:
\be
\D_{\D_0=4}=2R+\frac{2}{R}+O(R^{-2});
\ee
\item level-6 descendants, $\D_0=6$:
\be
\D_{\D_0=6}=2R+2+\frac{2}{R}+O(R^{-2}).
\ee
\end{itemize}

Table 1 of \texttt{1102.1209} contains in particular the following useful representatives:
\begin{itemize}
\item $\D_0=2$: the singlet $[0,0,0]_{(0,0)}$.
\item $\D_0=4$: the familiar descendants $[0,2,0]_{(1,1)}$ and $[2,0,2]_{(0,0)}$.
\item $\D_0=6$: the spin-4 singlet $[0,0,0]_{(2,2)}$ and the sphere highest-weight state
$[0,4,0]_{(0,0)}$.
\end{itemize}

The basic lesson is that the $\D_0=4$ states are best if one wants the simplest expected answer
$\D=2R+\frac{2}{R}+\cdots$, whereas the $\D_0=6$ states are easier to isolate in the
$SO(1,4)\times SO(5)$ language.

\section{Why the First Constant Correction Often Disappears}

This is the main structural issue in the SFT setup.

Suppose one starts from a pure NSNS trial state $\varphi^{(0)}_{NSNS}$. Then at first order,
\be
Q\varphi^{(1)}=-[\psi^{(1)}_{RR}\otimes \varphi^{(0)}_{NSNS}]',
\ee
so the first-order correction lives in the RR sector. Conversely, if one starts from a pure RR seed,
the first-order correction is driven into the NSNS sector.

Schematically, the first correction to the scaling dimension can only come from matrix elements of the form
\be
\delta \D^{(1)}\sim
\langle \varphi^{(0)}|H^{(1)}|\varphi^{(0)}\rangle
+\langle \varphi^{(0)}|H^{(0)}|\varphi^{(1)}\rangle,
\ee
with $H^{(1)}$ in the RR sector. For a pure NSNS seed, the diagonal term
$\langle \varphi^{(0)}_{NSNS}|H^{(1)}_{RR}|\varphi^{(0)}_{NSNS}\rangle$ is absent by sector/parity.
Thus one only gets the constant correction if the first-order source produces a nonzero opposite-sector
admixture with the same manifest quantum numbers.

This is where the simple candidates become deceptive:
\begin{itemize}
\item for some trial states the source $[\psi^{(1)}_{RR}\otimes \varphi^{(0)}]'$ vanishes outright by gamma-matrix kinematics;
\item for other trial states the source is nonzero, but there is no matching opposite-sector state with the same
manifest $SO(1,4)\times SO(5)$ labels and worldsheet parity to mix with;
\item a universal redefinition $\lambda^{1/4}=R+a_0+a_1/R+\cdots$ can reshuffle the constant and $1/R$ coefficients,
but it cannot by itself explain a state-by-state absence of the correction structure expected for a fixed supermultiplet.
\end{itemize}

So if a simple one-component ansatz shows no first constant correction, the right conclusion is not
that Konishi has $b_1=0$. The right conclusion is that the ansatz is too small, or too kinematically special,
to capture the needed opposite-sector mixing.

\section{Negative Control: The AdS Spin-4 Singlet}

The cleanest representative for the linearized equations is the level-6 AdS highest-weight state
$[0,0,0]_{(2,2)}$, i.e. the $S=4$ singlet. Introduce
\be
X^\pm = X^1 \pm i X^2,
\ee
and take the zeroth-order NSNS state
\be
\varphi^{(0)}_{S=4}
=a^{(0)}(X^+) \,
c\bar c\, e^{-\phi}\psi^+\partial X^+\,
e^{-\bar\phi}\bar\psi^+\bar\partial X^+.
\label{S4ansatz}
\ee
In the notation used later for the $\D_0=4$ descendants, this state has the trivial internal profile
\be
Y_{(0)}=1.
\ee
That is the basic difference between the $\D_0=6$ spin-4 singlet and the $\D_0=4$ state built on the same
$\psi^+\partial X^+$ oscillator core: for the singlet, all nontrivial quantum numbers are carried by the
polarization tensor in the $AdS$ directions, with no extra $S^5$ harmonic turned on.

This is attractive for several reasons.
\begin{itemize}
\item It is a single highest-weight component under the manifest $SO(1,4)$.
\item It is an $R$-singlet, so it is closer to ``Konishi proper'' than a pure sphere excitation.
\item Its zeroth-order BRST conditions are simple and already worked out in \texttt{draftKonishi.tex}:
\be
\left(1-\frac{1}{2}\Box\right)a^{(0)}=0,\qquad \partial^+ a^{(0)}=0.
\label{S4BRST}
\ee
\item The second-order source generated by $(h^{(2)},\phi^{(2)},A^{(2)})$ is nontrivial, so it is a good
warm-up problem for the actual linearized EOM technology.
\end{itemize}

However, this state is also the sharpest example of the first-order obstruction.
The first-order bracket computed in \texttt{draftKonishi.tex} is proportional to
\be
(\gamma^+F^{(1)}\gamma^+)^{\alpha\beta},
\ee
and vanishes because
\be
\gamma^+F^{(1)}\gamma^+ =0.
\ee
So the same feature that makes the mode easy to isolate also removes the obvious constant
correction in the one-mode truncation.

If this state is indeed a member of the Konishi multiplet, the expected answer is still
\be
\D_{S=4}=2R+2+\frac{2}{R}+O(R^{-2}),
\ee
because here $\D_0=6$. The point is that the first correction, namely the constant $+2$, is not visible
from the naive first-order diagonal analysis of the simple ansatz \eqref{S4ansatz}. The $2/R$ piece is then
a further-order question on top of that.

This is not a true contradiction. The representation-theoretic statement from Table 1 is only that there is
a unique \emph{interacting} Konishi state with labels $[0,0,0]_{(2,2)}$. It does \emph{not} imply that, in the
flat-space large-$R$ expansion, this state is already represented by the single bare NSNS vertex
\eqref{S4ansatz} with no additional homogeneous admixtures. Indeed, the first-order equation is
\be
Q\varphi^{(1)}_{S=4}=-[\varphi^{(0)}_{S=4}\otimes\psi^{(1)}]',
\ee
and for the simple highest-weight seed the right-hand side vanishes. But this only implies
\be
Q\varphi^{(1)}_{S=4}=0,
\ee
not $\varphi^{(1)}_{S=4}=0$. So first order leaves a nontrivial homogeneous sector, and the true Konishi
eigenstate may already require additional BRST-closed NSNS and/or RR components with the same overall
quantum numbers.

Equivalently, one should think of this as a degenerate-perturbation-theory problem on the full mass-level-1
state space, not as a one-by-one matrix element computation. The vanishing of
\be
\gamma^+F^{(1)}\gamma^+=0
\ee
only shows that the \emph{diagonal} first-order coupling of the top null-polarization component vanishes.
It does not prove that the whole first-order perturbation matrix vanishes in the enlarged level-1 sector.

There is also direct evidence in this direction from the higher-order equations already analyzed in
\texttt{draftKonishi.tex}: the spin-4 seed does not remain a closed one-field subsector. At second order it
sources additional $X^-$-dependent and transverse companions, so the single vertex \eqref{S4ansatz} was never
the full perturbative state to begin with. The correct interpretation is therefore:
\begin{itemize}
\item the spin-4 singlet is uniquely identified as a Konishi descendant at the representation-theory level;
\item the simple NSNS seed \eqref{S4ansatz} is only one convenient leading component of that state;
\item the apparent first-order contradiction is an artifact of over-truncating the degenerate level-1 problem.
\end{itemize}

\section{Secondary Control: The Sphere Highest-Weight Mode}

A similarly simple control problem is the sphere highest-weight state at $\D_0=6$.
Choose a complex polarization
\be
U^\pm = X^5 \pm i X^6,
\ee
and consider
\be
\varphi^{(0)}_{J=4}
=Y_{(4)}(U)\, c\bar c\, e^{-\phi}\psi^u\partial U^+\,
e^{-\bar\phi}\bar\psi^u\bar\partial U^+,
\ee
with
\be
Y_{(4)}(U)=\frac{(U^+)^4}{R^4}
=\frac{(X^5+iX^6)^4}{R^4},
\ee
where $\psi^u$ denotes the fermion in the same complex sphere direction.
This is the quartic highest-weight harmonic appropriate to the $[0,4,0]$ representation.
Equivalently one can think of it as the rank-4 symmetric traceless profile $(n\cdot X)^4/R^4$ with the
same null polarization $n^i=(1,i,0,0,0)$ used for $Y_{(2)}$.
If one prefers the notation used earlier in the draft, one may simply regard the old
$b^{(0)}(U^+)$ as shorthand for this quartic highest-weight profile.

This state is also straightforward in the reduced-symmetry frame: it is a single highest-weight
component under the manifest $SO(5)$. It should have the same strong-coupling form
\be
\D_{J=4}=2R+2+\frac{2}{R}+O(R^{-2}).
\ee

I do not currently see a reason to prefer it over the AdS spin-4 state for the first pass.
It is useful mainly as a control: if both simple highest-weight states fail to show the first-order
constant correction, then the issue is probably structural rather than specific to the AdS polarization.

\section{Better \texorpdfstring{$\D_0=4$}{Delta0=4} Candidates}

If the criterion is ``easiest from the SFT point of view,'' the natural place to look is the
$\D_0=4$ part of Table 1 rather than the $\D_0=6$ spin-4 singlet.
At $\D_0=4$ the first correction vanishes,
\be
\D_{\D_0=4}=2R+\frac{2}{R}+O(R^{-2}),
\ee
so one can focus directly on the nontrivial coefficient $b_1$.
The important candidates are
\be
[2,0,2]_{(0,0)},\qquad [0,2,0]_{(1,1)},\qquad [1,1,1]_{(\frac{1}{2},\frac{1}{2})},
\ee
with the caveat that the last two families come with multiplicities in Table 1.

To make the relevant zeroth-order quantum numbers visible already in the vertex operator, it is useful
to pick one complex $AdS$ plane and two complex sphere planes,
\be
X^\pm = X^1 \pm i X^2,\qquad
U^\pm = X^5 \pm i X^6,\qquad
V^\pm = X^7 \pm i X^8,
\ee
and to denote by $Y_{(2)}(U)$ a charge-two highest-weight sphere profile whose leading harmonic is
\be
Y_{(2)}(U)=\frac{(U^+)^2}{R^2}
=\frac{(X^5+iX^6)^2}{R^2}.
\ee
In the reduced-symmetry language this profile is the simplest way to keep track of the $S^5$
highest-weight data without committing to a full $SO(6)$-covariant harmonic basis.
Equivalently one can write
\be
Y_{(2)}(X)=\frac{1}{R^2}C_{ij}X^iX^j,
\ee
with $C_{ij}$ a symmetric traceless highest-weight tensor. For the present choice,
\be
C_{55}=1,\qquad C_{66}=-1,\qquad C_{56}=C_{65}=i,
\ee
with all other components vanishing. No extra trace subtraction is needed, because this is just the
null-polarization form $(n\cdot X)^2/R^2$ with
\be
n^i=(1,i,0,0,0),\qquad n_i n^i=0.
\ee
Thus the quadratic polynomial is already traceless and transforms in the desired $[0,2,0]$
highest-weight representation.

\subsection{Default target: \texorpdfstring{$[2,0,2]_{(0,0)}$}{[2,0,2](0,0)}}

This is the $su(2)$ Konishi descendant, explicitly identified in \texttt{1102.1209} through the operator
$\Tr([\Phi_1,\Phi_2]^2)$.
It is the best default candidate for the present SFT setup for three reasons.
\begin{itemize}
\item It is a bona fide Konishi descendant, not merely a candidate dual.
\item At $\D_0=4$ it appears as a single $[2,0,2]_{(0,0)}$ state in Table 1, unlike the
$[0,2,0]$ and $[1,1,1]$ families which come with multiplicities.
\item It is spinless in the $AdS_5$ directions, so its nontrivial quantum numbers live entirely in
$S^5$, which is closer to the manifest subgroup structure of the current SFT variables.
\end{itemize}

The natural zeroth-order NSNS representative in the $(-1,-1)$ picture is
\be
\varphi^{(0)}_{[2,0,2]}
=Y_{(2)}(U)\, c\bar c\,
e^{-\phi}\psi^{V^+}\partial V^+\,
e^{-\bar\phi}\bar\psi^{V^+}\bar\partial V^+.
\label{202vertex}
\ee
Its zeroth-order quantum numbers are easy to read off:
\begin{itemize}
\item there is no $AdS$ polarization, so the state is spinless in the chosen $AdS$ plane;
\item the profile $Y_{(2)}(U)$ carries a highest-weight sphere charge two in the $U$-plane;
\item the left- and right-moving oscillators in the $V$-plane contribute the second internal excitation.
\end{itemize}
So \eqref{202vertex} is the reduced-symmetry avatar of the $[2,0,2]_{(0,0)}$ descendant.

At zeroth order one expects the same type of BRST conditions as for the spin-4 test state
\eqref{S4ansatz}: the flat-space level-1 mass-shell condition together with transversality in the
polarized directions. For the particular highest-weight choice above, the transversality part is automatic
at the candidate level because the profile depends only on $U^+$ and is independent of the polarized
$V$-direction. In other words, \eqref{202vertex} is already the minimal highest-weight/traceless NSNS seed;
the real question is not zeroth-order identification but what RR partners it sources at first order.

From the SFT point of view, the natural hope is that one can start from a relatively small NSNS seed
built from two internal complex polarizations and only enlarge the ansatz if the first-order RR source
forces it.

The consistency checks for this candidate are:
\begin{itemize}
\item verify that $[\psi^{(1)}_{RR}\otimes \varphi^{(0)}_{[2,0,2]}]'$ does not vanish by a trivial gamma-matrix identity;
\item determine whether the first-order source lands in the $[1,1,1]_{(\frac{1}{2},\frac{1}{2})}$ RR family;
\item check that the second-order system preserves the traceless $[2,0,2]$ structure rather than collapsing onto singlet or trace components.
\end{itemize}

\subsection{Strong alternative: \texorpdfstring{$[0,2,0]_{(1,1)}$}{[0,2,0](1,1)}}

This is the standard $sl(2)$ Konishi descendant, explicitly identified in \texttt{1102.1209} through
$\Tr(\Phi_1 D^2\Phi_1)$ and matched there to $b_1=2$.
It has the strongest external evidence as a Konishi state, but it is slightly less attractive as a first SFT target.
\begin{itemize}
\item It mixes AdS spin and $S^5$ charge, so it is less naturally packaged in the manifest subgroup language.
\item The $[0,2,0]$ family at $\D_0=4$ is not isolated to a single state in Table 1.
\item One has to make sure that the ansatz isolates the $(1,1)$ component rather than drifting into one of the scalar $[0,2,0]_{(0,0)}$ copies.
\end{itemize}

Its main consistency checks are:
\begin{itemize}
\item verify that the $J=2$ structure can be represented economically in the embedding-coordinate variables;
\item check whether the first-order RR source is less degenerate than for the spin-4 singlet;
\item confirm that the resulting system does not mix uncontrollably with the other $[0,2,0]$ states at the same level.
\end{itemize}

The corresponding zeroth-order NSNS representative is
\be
\varphi^{(0)}_{[0,2,0]}
=Y_{(2)}(U)\, c\bar c\,
e^{-\phi}\psi^+\partial X^+\,
e^{-\bar\phi}\bar\psi^+\bar\partial X^+.
\label{020vertex}
\ee
Its quantum numbers are again immediate from the operator itself:
\begin{itemize}
\item the oscillator part is the same $(1,1)$ highest-weight AdS excitation as in the spin-4 control,
but with only one $X^+$ derivative on each side;
\item the sphere profile contributes the charge-two internal highest weight;
\item together these give the reduced-symmetry version of the $[0,2,0]_{(1,1)}$ Konishi descendant.
\end{itemize}

It is important that this is \emph{not} a different choice of oscillator core from the earlier
suggestion $c\bar c\,e^{-\phi}\psi^+\partial X^+\,e^{-\bar\phi}\bar\psi^+\bar\partial X^+$.
The AdS part is exactly the same highest-weight $(1,1)$ excitation. The difference is the extra factor
$Y_{(2)}(U)$.
Without that sphere profile one is only specifying the $AdS$ quantum numbers, which is not enough to
distinguish the $\D_0=4$ descendant from the $\D_0=6$ spin-4 singlet in the reduced-symmetry language.
The role of $Y_{(2)}(U)$ is precisely to supply the internal charge-two data that moves the state into the
$[0,2,0]$ family.

As for \eqref{S4ansatz}, the zeroth-order BRST conditions should reduce to the mass-shell equation for
the profile together with transversality. For \eqref{020vertex} the transversality condition is again
automatic at the level of this highest-weight seed, since $Y_{(2)}(U)$ is independent of the conjugate
$AdS$ direction. So \eqref{020vertex} is the natural one-component starting point for the $sl(2)$
descendant, with the understanding that the full first-order solution may force RR companions.

\subsection{Likely RR partners: \texorpdfstring{$[1,1,1]_{(\frac{1}{2},\frac{1}{2})}$}{[1,1,1](1/2,1/2)}}

These are not my preferred standalone starting points, but they are important because the odd-order RR
background naturally drives a first-order correction of an NSNS seed into RR-type states.
At $\D_0=4$ Table 1 contains two such $[1,1,1]_{(\frac{1}{2},\frac{1}{2})}$ states.

From the SFT point of view, these should be treated as probable mixing partners rather than primary seeds.
The key check is simply whether a chosen NSNS candidate sources this RR family already at first order.
If yes, the minimal sensible ansatz should include them from the start instead of insisting on a pure NSNS truncation.

\subsection{Ambiguous scalar option: \texorpdfstring{$[0,2,0]_{(0,0)}$}{[0,2,0](0,0)}}

This scalar is worth remembering because \texttt{1102.1209} mentions it as a candidate dual to a pulsating
string with $N=2$ and $J=2$. It looks superficially easy from the SFT point of view because it is spinless in $AdS$.
I do not think it should be the default target, because the identification is weaker than for
$[2,0,2]_{(0,0)}$ or $[0,2,0]_{(1,1)}$, and the $[0,2,0]$ family is already degenerate at $\D_0=4$.

\subsection{Why Their First-Order EOM Match the Expected Dimension}

For either of the zeroth-order NSNS seeds \eqref{202vertex} or \eqref{020vertex}, the first correction obeys
\be
Q\varphi^{(1)}=-[\psi^{(1)}_{RR}\otimes \varphi^{(0)}_{\D_0=4}]'.
\label{FirstOrderDelta4}
\ee
The right-hand side is RR, so the first correction is generically RR-valued. This is exactly what one
should expect from the alternating RR/NSNS structure of the background.

Crucially, this does \emph{not} conflict with the Konishi dimension formula for the $\D_0=4$ descendants,
because for these states the first constant correction is supposed to vanish:
\be
\D_{\D_0=4}=2R+\frac{2}{R}+O(R^{-2}).
\ee
So the first-order equation \eqref{FirstOrderDelta4} is not required to generate a diagonal mass shift.
Its job is only to determine the opposite-sector admixture demanded by the RR background.

There are then two logically allowed outcomes, and both are consistent with the expected dimension.
\begin{itemize}
\item If $[\psi^{(1)}_{RR}\otimes \varphi^{(0)}_{\D_0=4}]'\neq 0$, then the true first-order eigenstate is
not purely NSNS: it contains an RR piece, most naturally in the $[1,1,1]_{(\frac{1}{2},\frac{1}{2})}$ family.
This is compatible with $b_0=0$ because no constant shift is required.
\item If $[\psi^{(1)}_{RR}\otimes \varphi^{(0)}_{\D_0=4}]'=0$, that is also compatible with $b_0=0$.
It simply means that the first nontrivial dynamical test moves to the next order, where one should recover
the universal Konishi coefficient $b_1=2$.
\end{itemize}

This is precisely why the $\D_0=4$ candidates are better behaved than the $\D_0=6$ controls.
For the spin-4 singlet and the $J=4$ sphere mode one expects
\be
\D_{\D_0=6}=2R+2+\frac{2}{R}+O(R^{-2}),
\ee
so a vanishing first-order source is then a genuine warning sign: the simple pure NSNS seed cannot by
itself reproduce the required constant $+2$, and one probably needs a mixed zeroth-order ansatz already at
the start.

\subsection{Explicit First-Order Bracket Check for the \texorpdfstring{$\D_0=4$}{Delta0=4} Seeds}
\label{sec:delta4_bracket}

We can now perform the gamma-matrix check flagged in \S6.1 for both $\D_0=4$ candidates.
The outcome will make clear why the $\D_0=4$ seeds are ``better behaved'' in a more precise sense
than the discussion in \S6.5 alone captures.

\paragraph{The $[2,0,2]_{(0,0)}$ seed.}
For the vertex \eqref{202vertex},
\be
\varphi^{(0)}_{[2,0,2]}
=Y_{(2)}(U)\, c\bar c\,
e^{-\phi}\psi^{V^+}\partial V^+\,
e^{-\bar\phi}\bar\psi^{V^+}\bar\partial V^+,
\ee
both worldsheet sectors carry the same complexified null direction $V^+=X^7+iX^8$ in $S^5$.
The relevant gamma-matrix bilinear is
\be
\gamma^{V^+}F^{(1)}\gamma^{V^+}
=\frac{i}{R}\,\gamma^{V^+}\gamma^{01234}\gamma^{V^+}.
\ee
Since $V^+$ is an $S^5$ direction, $\gamma^{V^+}$ anticommutes with $\gamma^{01234}$,
so one can pull the left factor through:
$\gamma^{V^+}\gamma^{01234}\gamma^{V^+}=-\gamma^{01234}(\gamma^{V^+})^2$.
But
\be
(\gamma^{V^+})^2=(\gamma^7+i\gamma^8)^2
=(\gamma^7)^2+i\{\gamma^7,\gamma^8\}+(i\gamma^8)^2
=-1+0+1=0,
\ee
where the signs come from the Euclidean signature of the $S^5$ directions ($(\gamma^i)^2=-1$).

\emph{The first-order bracket vanishes for the $[2,0,2]_{(0,0)}$ highest-weight seed,
by the same null-polarisation mechanism as for the spin-4 singlet.}

\paragraph{The $[0,2,0]_{(1,1)}$ seed.}
For the vertex \eqref{020vertex},
\be
\varphi^{(0)}_{[0,2,0]}
=Y_{(2)}(U)\, c\bar c\,
e^{-\phi}\psi^+\partial X^+\,
e^{-\bar\phi}\bar\psi^+\bar\partial X^+,
\ee
both sectors carry $X^+=X^1+iX^2$ in $AdS$.
The bilinear is $\gamma^+\gamma^{01234}\gamma^+=\gamma^{01234}(\gamma^+)^2=0$
(with $(\gamma^+)^2=0$ as in \S4). So the bracket vanishes here too.

\paragraph{Interpretation.}
The vanishing of both first-order brackets is consistent with $b_0=0$ for $\D_0=4$
and changes nothing about the expected scaling dimension
$\D=2R+2/R+O(R^{-2})$.
The first nontrivial dynamical coefficient is $b_1=2$, which enters at second order.
Therefore the $\D_0=4$ seeds are well suited for a \emph{second-order} computation
of the Konishi anomalous dimension.

\paragraph{Why the first-order source vanishes for all highest-weight seeds.}
The common origin of the vanishing in all cases above is the use of a complexified
null polarisation on both worldsheet sectors. For any vector $n^\mu$ with $n_\mu n^\mu=0$
(whether in $AdS$ or in $S^5$), the combination
$\gamma^{n}=n_\mu\gamma^\mu$ satisfies $(\gamma^n)^2=0$,
so the bilinear $\gamma^n\gamma^{01234}\gamma^n$ always vanishes.
Since the natural highest-weight seeds always choose the same null direction for both the
left- and right-moving oscillators, the first-order source vanishes for all of them.

There are two ways to proceed:
\begin{enumerate}
\item \textbf{Use a state whose SO(5) quantum numbers force different directions on the two
worldsheet sectors,} so that the first-order source is nonzero. This is the mechanism exploited
by $\langle 4,0\rangle_{(1,1)}$ in \S\ref{sec:202_11}: the Weyl-tensor structure
$T_{[ij],[kl]}$ requires the pairs $(ij)$ and $(kl)$ to be in \emph{different} planes,
so $\gamma^i\gamma^{01234}\gamma^k\neq 0$.
\item \textbf{Accept the vanishing first-order source and go directly to second order.}
This is what the $\D_0=4$ candidates allow, since $b_0=0$ is expected and
no information is lost.
\end{enumerate}

For the second-order computation itself, it is worth noting that the second-order bracket
$[\psi^{(2)}_{NSNS}\otimes\varphi^{(0)}]'$ does \emph{not} have a vanishing source
(it involves bosonic profiles $h_{\mu\nu}^{(2)}$, $\phi^{(2)}$, $A_\mu^{(2)}$
contracted with the seed, as already worked out in detail for the spin-4 singlet in
\texttt{draftKonishi.tex}).

\paragraph{Summary for the $\D_0=4$ strategy.}
The $\D_0=4$ seeds \eqref{202vertex} and \eqref{020vertex} cannot be uniquely identified by
$SO(1,4)\times SO(5)$ quantum numbers (their manifest labels are shared by states at
$\D_0=4,5,6,7,8$). However, if one accepts the identification on other grounds
(operator matching, Table-1 uniqueness within each $[p,q,r]$, or descent from a uniquely
identified $\D_0=6$ state), then $b_0=0$ makes them the cleanest setting for extracting $b_1=2$
at second order.

The computation for $[2,0,2]_{(0,0)}$ is the more economical one: both oscillators are
internal, the profile is a simple quadratic harmonic, and the state is spinless in $AdS$.
For $[0,2,0]_{(1,1)}$, the AdS spin forces additional bookkeeping in the second-order
brackets, but the external evidence for the identification is strongest.

\subsection{How Far Manifest \texorpdfstring{$SO(5)$}{SO(5)} Really Gets Us}

The manifest $SO(5)$ is sufficient to write down concrete highest-weight seeds, but it is not by itself
enough to certify the full $SO(6)$ representation, let alone Konishi membership.

For the concrete candidates above, the practical situation is as follows.
\begin{itemize}
\item For the spin-4 control \eqref{S4ansatz}, the internal profile is just $Y_{(0)}=1$. Here the loss of
full $SO(6)$ is not the main issue: the state is already an internal singlet. The real problem is dynamical,
namely the first-order obstruction.
\item For the $sl(2)$ candidate \eqref{020vertex}, the explicit profile
$Y_{(2)}=(X^5+iX^6)^2/R^2$ is a highest-weight component of the rank-2 traceless tensor of the manifest
$SO(5)$. This is enough to say that one is looking at the correct \emph{type} of internal excitation.
But under the standard branching $SO(6)\to SO(5)$, the full $[0,2,0]$ decomposes as
\be
[0,2,0]\downarrow SO(5)=14\oplus 5\oplus 1.
\ee
So the manifest description only picks one branch of the full $SO(6)$ multiplet. It does not by itself
reconstruct how the hidden generators mixing the fifth-sphere direction with the non-manifest sixth one
complete that branch into the full $[0,2,0]$ representation.
\item The same issue is even clearer for the quartic control profile
$Y_{(4)}=(X^5+iX^6)^4/R^4$. As an $SO(5)$ tensor it is the highest-weight component of the rank-4
traceless symmetric tensor, but the full $SO(6)$ representation branches as
\be
[0,4,0]\downarrow SO(5)=55\oplus 30\oplus 14\oplus 5\oplus 1.
\ee
Again, the manifest variables isolate a very useful top component, but not the whole $SO(6)$ multiplet.
\item For the preferred $su(2)$ candidate \eqref{202vertex}, manifest $SO(5)$ gets us somewhat farther in
practice because both the profile $Y_{(2)}(U)$ and the oscillator polarization
$\psi^{V^+}\partial V^+$ live in the internal directions. So one can build a small seed with no $AdS$ spin
and with very explicit internal charges. This makes it a good \emph{working ansatz}. But it is still not a
proof of the full $[2,0,2]$ assignment, because the missing $SO(6)/SO(5)$ generators and the first-order
RR mixing can relate this seed to other states with the same manifest $SO(5)$ labels.
\end{itemize}

The right conclusion is therefore intermediate.
\begin{itemize}
\item Manifest $SO(5)$ is sufficient to choose economical zeroth-order trial vertices and to separate the
main qualitative options: singlet, rank-2 internal tensor, rank-4 internal tensor, mixed AdS/internal state,
etc.
\item Manifest $SO(5)$ is not sufficient to pin down ``this is the Konishi state'' in a representation-theoretic
sense. To establish that, one still needs the full linearized EOM together with the non-manifest symmetry
completion, or equivalently enough information to see that the chosen seed sits in the correct
$PSU(2,2|4)$ long multiplet.
\end{itemize}

\subsection{Explicit Table-1 Collision Check}
\label{app:table1scan}

One can make the previous statement completely explicit by scanning all 17 rows of Table 1 of
\texttt{1102.1209}. I branch every $SO(6)$ irrep down to manifest $SO(5)$ and denote the resulting
$SO(5)$ irreps by $\langle a,b\rangle$ in the $C_2\simeq B_2$ convention, with
\be
\langle 1,0\rangle = {\bf 4},\qquad \langle 0,1\rangle = {\bf 5}.
\ee
The full branching list used in the scan is collected in Appendix~\ref{app:branching}, and Table 1 itself
is reproduced in Appendix~\ref{app:table1}.

The important caveat is that Table 1 is \emph{not} already organized in manifest
$SO(1,4)\times SO(5)$ multiplets. On the spacetime side it gives the boundary
$SO(4)\subset SO(2,4)$ labels $(s_L,s_R)$, not a bulk $SO(1,4)$ tensor label.
So the direct all-state question
\be
\text{``which full }SO(1,4)\times SO(5)\text{ irrep does this Table-1 state become?''}
\ee
cannot actually be answered from Table 1 alone.

One can see this explicitly by trying to lift the displayed $SO(4)$ content to finite-dimensional
$SO(1,4)$ representations. The simplest obstruction already appears at $\D_0=3$:
\be
[0,1,0]_{(0,1)+(1,0)}.
\ee
There is no finite-dimensional $SO(1,4)$ irrep whose $SO(4)$ branching is just $(0,1)+(1,0)$.
The minimal candidate would be $\langle 2,0\rangle$, but
\be
\langle 2,0\rangle \downarrow SO(4)=(0,1)+\Big(\frac12,\frac12\Big)+(1,0),
\ee
so the $(\frac12,\frac12)$ piece is missing. The same issue persists after first branching the
internal side to $SO(5)$: for instance, the full $\langle 0,2\rangle$ block at $\D_0=4$ has total
$SO(4)$ content
\be
3(0,0)+(1,1),
\ee
whereas a local symmetric traceless $SO(1,4)$ tensor $\langle 0,2\rangle$ would branch as
\be
\langle 0,2\rangle \downarrow SO(4)=(0,0)+\Big(\frac12,\frac12\Big)+(1,1).
\ee
Thus Table 1 does \emph{not} supply a full manifest $SO(1,4)$ label for every state.

What one \emph{can} do explicitly is a lower-resolution scan that keeps the full $SO(5)$ branching
on the internal side and the displayed $(s_L,s_R)$ data on the spacetime side. Concretely, the check
uses the pair
\be
\big(\langle a,b\rangle,\,(s_L,s_R)\big)
\ee
as a coarse label. If collisions already occur at this level, then Table 1 certainly does not give
enough information to distinguish those states in a manifest $SO(1,4)\times SO(5)$ language.

The outcome of this coarse scan is still stark. After branching the full Table 1 in this way, one finds
\be
109
\ee
distinct manifest labels $\big(\langle a,b\rangle,(s_L,s_R)\big)$.
Of these,
\be
97
\ee
occur at more than one value of $\D_0$.
Only
\be
12
\ee
manifest labels are unique to a single value of $\D_0$, and only
\be
11
\ee
actually isolate a single Table-1 state. All 11 of these fully unique labels sit at $\D_0=6$.

So the correct conclusion is:
\begin{itemize}
\item Table 1 by itself does not determine the full $SO(1,4)\times SO(5)$ branching of all Konishi states.
\item Even at the coarser level accessible directly from Table 1, it is very common for states with different
$\D_0$ to look identical.
\end{itemize}

Some concrete collisions are especially relevant for the present SFT discussion.
\begin{itemize}
\item The scalar singlet label
\be
\langle 0,0\rangle_{(0,0)}
\ee
appears at
\be
\D_0=2,4,6,8,10,
\ee
coming from the $SO(6)$ representations $[0,0,0]$, $[0,2,0]$ and $[0,4,0]$.
\item The manifest label relevant to the simple scalar $su(2)$-type internal ansatz,
\be
\langle 0,2\rangle_{(0,0)},
\ee
appears at
\be
\D_0=4,5,6,7,8,
\ee
coming from $[0,2,0]$, $[1,1,1]$, $[1,2,1]$, $[0,4,0]$ and $[2,0,2]$.
So a rank-2 internal scalar profile by itself is far from enough to identify the $\D_0=4$ Konishi
descendant.
\item The manifest label relevant to the $sl(2)$ descendant,
\be
\langle 0,2\rangle_{(1,1)},
\ee
also appears at
\be
\D_0=4,5,6,7,8,
\ee
coming from $[0,2,0]$, $[1,1,1]$ and $[2,0,2]$.
\item By contrast, the AdS spin-4 singlet is special: its manifest label
\be
\langle 0,0\rangle_{(2,2)}
\ee
occurs only once in the whole table, namely for the $\D_0=6$ state $[0,0,0]_{(2,2)}$.
This explains why it was so easy to isolate in the reduced-symmetry variables.
\item The sphere highest-weight control also has a unique top branch:
\be
\langle 0,4\rangle_{(0,0)}
\ee
appears only once, for $[0,4,0]_{(0,0)}$ at $\D_0=6$.
\item A similar statement holds for the top branch
\be
\langle 4,0\rangle_{(1,1)},
\ee
which appears only once, for the $\D_0=6$ state $[2,0,2]_{(1,1)}$.
\end{itemize}

The practical lesson is therefore very sharp.
\begin{itemize}
\item If one works only with the data directly visible in Table 1, then most states in the Konishi table are
heavily degenerate across different values of $\D_0$.
\item A genuine $SO(1,4)$-refined scan requires extra bulk input beyond Table 1, i.e. a dictionary from the
Table-1 conformal data to actual $AdS_5$ field representations.
\item The apparently simple $\D_0=4$ candidates are \emph{not} protected against this ambiguity: their low
rank manifest $SO(5)$ branches are shared by many other states.
\item The reason the spin-4 singlet was such a clean diagnostic was precisely that it sits in one of the very
few manifest labels that remain unique after branching.
\end{itemize}

\section{A Useful Caution from RR Seeds}

One could try to avoid the previous problem by starting from an RR representative instead of an NSNS one.
But that does not solve the issue automatically. For a pure RR seed, the first-order correction must lie in
the NSNS sector, and one again needs an NSNS state with the same manifest quantum numbers.

The example already mentioned in \texttt{draftKonishi.tex} is a level-3 RR descendant with $\D_0=5$:
it has the same flat-space mass as the simple NSNS Konishi candidate, but there is no obvious NSNS state
of the required spin to mix with at first order. So pure RR seeds can fail for the same structural reason.

In other words, the problem is not ``RR bad, NSNS good'' or vice versa. The real issue is whether the
reduced-symmetry sector one is working in contains enough states to support the odd-order RR/even-order NSNS
mixing forced by the background.

\section{Working Recommendation}

\subsection{Where \texorpdfstring{$b_0$}{b0} comes from}

Since $H^{(1)}=0$ (the Hamiltonian has no first-order correction), the first-order
perturbation matrix is trivially zero and $b_0$ cannot arise at first order.
Moreover, $F^{(1)}$ is constant, so the first-order linearized EOM preserves the harmonic
level $\ell$ of each state (\S\ref{sec:b0resolution}).

Both $b_0=\D_0-4$ and $b_1$ must therefore emerge from the \emph{second-order} Hamiltonian
$H^{(2)}$, which involves the NSNS background $\psi^{(2)}$ with quadratic $X$-dependence.
The $X_iX_j$ part of $h^{(2)}_{\mu\nu}$ mixes harmonic levels $\ell$ and $\ell\pm 2$,
coupling the $\D_0=6$ ($\ell=0$), $\D_0=5$ ($\ell=1$), and $\D_0=4$ ($\ell=2$) states.

Diagonalising $H^{(2)}$ in this combined space should simultaneously produce $b_0=+2,+1,0$
for the three harmonic levels and determine $b_1$ for each. This is the computation already
underway in \texttt{draftKonishi.tex}.

\subsection{Recommended strategy: \texorpdfstring{$\D_0=4$}{Delta0=4} candidates}

The $\D_0=4$ candidates do not have this problem: $b_0=\D_0-4=0$, and the first nontrivial
coefficient is $b_1=2$ at second order.

Take $[2,0,2]_{(0,0)}$ at $\D_0=4$ as the default computational target.
\begin{itemize}
\item It is a bona fide Konishi descendant, unique in Table~1 at $\D_0=4$,
spinless in $AdS$, and naturally carried by internal quantum numbers.
\item The first-order source vanishes (null-polarisation identity, \S\ref{sec:delta4_bracket}),
which is consistent with $b_0=0$. No information is lost.
\item The first nontrivial coefficient $b_1=2$ enters at second order, where the NSNS bracket
$[\psi^{(2)}\otimes\varphi^{(0)}]'$ is nonvanishing and already worked out for the spin-4
singlet in \texttt{draftKonishi.tex}.
\item The vertex operator is simple: a quadratic scalar harmonic $Y_{(2)}(U)$ times internal
oscillators $\psi^{V^+}\partial V^+$ on each side.
\item The manifest label $\langle 0,2\rangle_{(0,0)}$ is \emph{not} unique: it appears at
$\D_0=4,5,6,7,8$. The identification rests on Table~1 uniqueness within the $[2,0,2]$
family, not on manifest quantum numbers alone.
\end{itemize}

Keep $[0,2,0]_{(1,1)}$ as the strongest alternative if external confidence in the identification
matters more than ansatz economy.

\section{The \texorpdfstring{$[2,0,2]_{(1,1)}$}{[2,0,2](1,1)} Candidate at \texorpdfstring{$\D_0=6$}{Delta0=6}}
\label{sec:202_11}

The analysis in \S\ref{app:table1scan} shows that only 11 manifest labels
$\big(\langle a,b\rangle,(s_L,s_R)\big)$ isolate a single Table-1 state, all at $\D_0=6$.
Among these, the simplest candidates that were checked --- the spin-4 singlet
$\langle 0,0\rangle_{(2,2)}$ and the sphere $J\!=\!4$ mode $\langle 0,4\rangle_{(0,0)}$
--- both have a vanishing first-order source, as discussed in \S4 and \S\ref{sec:delta4_bracket}.
This section identifies
\be
\langle 4,0\rangle_{(1,1)},\qquad\text{from }[2,0,2]_{(1,1)}\text{ at }\D_0=6,
\ee
as a candidate that is uniquely pinned down by manifest $SO(1,4)\times SO(5)$
quantum numbers \emph{and} has a nonvanishing first-order source.

\subsection{Representation theory: \texorpdfstring{$\langle 4,0\rangle$}{<4,0>} of \texorpdfstring{$SO(5)$}{SO(5)}}

In the $B_2$ convention used throughout ($\langle 1,0\rangle={\bf 4}$, $\langle 0,1\rangle={\bf 5}$),
a representation $\langle a,b\rangle$ of $SO(5)$ is tensorial (single-valued on $SO(5)$, not only
on $\mathrm{Spin}(5)$) if and only if $a$ is even.
This is because the first fundamental weight $\omega_1=(e_1+e_2)/2$
has half-integer coordinates in the orthonormal basis, so $\lambda=a\omega_1+b\omega_2$ has
integer coordinates iff $a$ is even.

Since $a=4$ is even, $\langle 4,0\rangle$ is a \emph{tensor} representation of $SO(5)$.
It can be realised entirely from tensor products of the vector representation $\langle 0,1\rangle={\bf 5}$.

Using the Weyl dimension formula for $B_2$,
\be
\dim\langle a,b\rangle = \frac{(a+1)(b+1)(a+2b+3)(a+b+2)}{6},
\ee
one finds $\dim\langle 4,0\rangle = 5\cdot 1\cdot 7\cdot 6/6 = 35$.

The explicit tensor product chain is
\ie
\langle 0,1\rangle\otimes\langle 0,1\rangle
&= \langle 0,0\rangle\oplus\langle 2,0\rangle\oplus\langle 0,2\rangle,\\
\langle 2,0\rangle\otimes\langle 2,0\rangle
&= \langle 4,0\rangle\oplus\langle 2,1\rangle\oplus\langle 0,2\rangle
\oplus\langle 2,0\rangle\oplus\langle 0,1\rangle\oplus\langle 0,0\rangle,
\fe
so $\langle 4,0\rangle$ sits inside $\mathrm{Sym}^2(\Lambda^2 V_5)$.
As a rank-4 tensor in $SO(5)$ vector indices, it has the symmetries of the Weyl tensor:
\be
T_{[ij],[kl]}=T_{[kl],[ij]},\qquad
T_{[ij],[kl]}+T_{[ik],[lj]}+T_{[il],[jk]}=0,\qquad
\delta^{ik}T_{[ij],[kl]}=0.
\ee

\paragraph{Scalar harmonics only give $\langle 0,b\rangle$.}
Polynomial harmonics on $S^4$ (homogeneous harmonic polynomials in $X^i$, $i=5,\ldots,9$)
furnish the symmetric traceless tensor representations $\langle 0,b\rangle$.
Since $\langle 4,0\rangle$ is \emph{not} of the form $\langle 0,b\rangle$, it cannot be
carried by a scalar profile alone.

\subsection{Vertex operator}

At the first massive level one has one holomorphic and one anti-holomorphic oscillator excitation.
Since the $\langle 4,0\rangle$ charge cannot sit in the profile, the worldsheet oscillators must carry
$S^5$ indices. The natural NSNS vertex operator in the $(-1,-1)$ picture is
\be\label{202_11vertex}
\varphi^{(0)}_{[2,0,2]_{(1,1)}}
= f(X^a)\,T_{[ij],[kl]}\,
c\bar c\,
e^{-\phi}\psi^i\partial X^j\,
e^{-\bar\phi}\bar\psi^k\bar\partial X^l,
\ee
where
\begin{itemize}
\item $i,j,k,l\in\{5,\ldots,9\}$ are $S^5$ directions;
\item $T_{[ij],[kl]}$ is the $\langle 4,0\rangle$ highest-weight projection
of the Weyl-tensor-type contraction (antisymmetric in each pair, symmetric under pair exchange,
with all traces removed);
\item $f(X^a)$ carries the boundary spin $(1,1)$ via its dependence on the $AdS_5$ coordinates
$X^0,\ldots,X^4$, together with the radial/mass-shell information appropriate to $\D_0=6$.
\end{itemize}

The highest-weight component of $T$ is obtained by choosing the null polarisations
\be
U^\pm = X^5\pm iX^6,\qquad V^\pm = X^7\pm iX^8,
\ee
and setting
\be
T_{[56],[78]}\sim 1,
\ee
i.e.\ the antisymmetric pair $(ij)=(56)$ and $(kl)=(78)$, both in the ``$+$'' direction
of their respective complex planes. The profile $f$ then depends only on $X^+\equiv X^1+iX^2$
and on the radial variable; concretely one expects $f\sim a(X^+)$ with $a$ satisfying the same
mass-shell and transversality conditions as before.

\subsection{Full mass-level-1 spectrum in the \texorpdfstring{$\langle 4,0\rangle_{(1,1)}$}{<4,0>(1,1)} sector}
\label{sec:spectrum}

To understand the first-order dynamics of this state we need the complete list of
mass-level-1 states with quantum numbers $\langle 4,0\rangle_{(1,1)}$, in both the NSNS and
RR sectors. This is a straightforward exercise in $SO(9)\to SO(4)\times SO(5)$
representation theory.

\paragraph{Physical states at mass level~1.}
The bosonic physical states of the Type~IIB closed superstring at the first massive level
form:
\ie
\text{NSNS:}&\quad ({\bf 44}\oplus{\bf 84})_L\otimes({\bf 44}\oplus{\bf 84})_R,\\
\text{RR:}&\quad {\bf 128}_L\otimes{\bf 128}_R,
\fe
where ${\bf 44}$ is the symmetric traceless tensor (rank 2), ${\bf 84}$ is the 3-form, and
${\bf 128}$ is the gamma-traceless vector-spinor, all of $SO(9)$.

\paragraph{Decomposition under $SO(4)\times SO(5)$.}
The vector ${\bf 9}$ of $SO(9)$ decomposes as
${\bf 9}\to\big((\tfrac12,\tfrac12),\langle 0,0\rangle\big)\oplus\big((0,0),\langle 0,1\rangle\big)$.

For the ${\bf 44}=\mathrm{Sym}^2_0({\bf 9})$ (traceless symmetric square):
\be\label{44decomp}
{\bf 44}\to
\big((0,0),\langle 0,0\rangle\big)
\oplus\big((1,1),\langle 0,0\rangle\big)
\oplus\big((\tfrac12,\tfrac12),\langle 0,1\rangle\big)
\oplus\big((0,0),\langle 0,2\rangle\big).
\ee
For the ${\bf 84}=\Lambda^3({\bf 9})$:
\ie\label{84decomp}
{\bf 84}\to{}&
\big((\tfrac12,\tfrac12),\langle 0,0\rangle\big)
\oplus\big((1,0)\oplus(0,1),\langle 0,1\rangle\big)
\nonumber\\
&\oplus\big((\tfrac12,\tfrac12),\langle 2,0\rangle\big)
\oplus\big((0,0),\langle 2,0\rangle\big).
\fe
Here $\langle 2,0\rangle$ in the ${\bf 84}$ arises from $\Lambda^1({\bf 4})\otimes\Lambda^2({\bf 5})$
(one $AdS$ and two $S^5$ indices, spin $(\tfrac12,\tfrac12)$) and from $\Lambda^3({\bf 5})$
(three $S^5$ indices, spin $(0,0)$).

For the ${\bf 128}$ (gamma-traceless vector-spinor): the spinor ${\bf 16}$ of $SO(9)$ decomposes as
${\bf 16}\to\big((\tfrac12,0),\langle 1,0\rangle\big)\oplus\big((0,\tfrac12),\langle 1,0\rangle\big)$,
and the ${\bf 128}\subset{\bf 9}\otimes{\bf 16}$ has
\ie\label{128decomp}
{\bf 128}\to{}&
\big((1,\tfrac12),\langle 1,0\rangle\big)
\oplus\big((\tfrac12,1),\langle 1,0\rangle\big)
\oplus\big((0,\tfrac12),\langle 1,0\rangle\big)
\oplus\big((\tfrac12,0),\langle 1,0\rangle\big)
\nonumber\\
&\oplus\big((\tfrac12,0),\langle 1,1\rangle\big)
\oplus\big((0,\tfrac12),\langle 1,1\rangle\big)
+\;\text{gamma-trace removal.}
\fe
The key point: the $\langle 1,1\rangle$ component of the ${\bf 128}$ only
carries $AdS$ spins $(\tfrac12,0)$ and $(0,\tfrac12)$.

\paragraph{NSNS states with $\langle 4,0\rangle_{(1,1)}$.}
From \eqref{44decomp} and \eqref{84decomp}, the $SO(5)$ representations
that can produce $\langle 4,0\rangle$ via the tensor product are
\be
\langle 2,0\rangle\otimes\langle 2,0\rangle\supset\langle 4,0\rangle,
\qquad
\langle 0,2\rangle\otimes\langle 2,0\rangle\supset\langle 4,0\rangle,
\qquad
\langle 0,2\rangle\otimes\langle 0,2\rangle\supset\langle 4,0\rangle.
\ee
Checking the $SO(4)$ spin content for each possibility:

\smallskip
\begin{center}
\begin{tabular}{|l|l|l|c|}
\hline
Sector & $SO(5)$ product & Spin product & $(1,1)$? \\
\hline
${\bf 84}\otimes{\bf 84}$ & $\langle 2,0\rangle\otimes\langle 2,0\rangle$ &
$(\tfrac12,\tfrac12)\otimes(\tfrac12,\tfrac12)\supset(1,1)$ & \textbf{Yes} \\
${\bf 84}\otimes{\bf 84}$ & $\langle 2,0\rangle\otimes\langle 2,0\rangle$ &
$(\tfrac12,\tfrac12)\otimes(0,0)=(\tfrac12,\tfrac12)$ & No \\
${\bf 84}\otimes{\bf 84}$ & $\langle 2,0\rangle\otimes\langle 2,0\rangle$ &
$(0,0)\otimes(0,0)=(0,0)$ & No \\
${\bf 44}\otimes{\bf 84}$ & $\langle 0,2\rangle\otimes\langle 2,0\rangle$ &
$(0,0)\otimes(\tfrac12,\tfrac12)=(\tfrac12,\tfrac12)$ & No \\
${\bf 84}\otimes{\bf 44}$ & $\langle 2,0\rangle\otimes\langle 0,2\rangle$ &
same as above & No \\
${\bf 44}\otimes{\bf 44}$ & $\langle 0,2\rangle\otimes\langle 0,2\rangle$ &
$(0,0)\otimes(0,0)=(0,0)$ & No \\
\hline
\end{tabular}
\end{center}
\smallskip

\noindent\textbf{NSNS total: 1 state}, from ${\bf 84}_L\otimes{\bf 84}_R$.

In terms of $SO(9)$ tensors, this state descends from the component $C_{a,[ij]}$
(one $AdS$ vector index, two antisymmetric $S^5$ indices) on each side. The
$\langle 4,0\rangle$ is assembled from $\langle 2,0\rangle_L\otimes\langle 2,0\rangle_R$,
and the spin $(1,1)$ from $(\tfrac12,\tfrac12)_L\otimes(\tfrac12,\tfrac12)_R$.

Note that in terms of vertex operators, the $C_{a,[ij]}$ component of the ${\bf 84}$ involves
both the two-oscillator structure $\psi^{[i}\partial X^{j]}$ and the three-fermion structure
$\psi^a\psi^i\psi^j$ (they mix under BRST). The vertex \eqref{202_11vertex} as written
includes only the former; the physical BRST-closed vertex will generically involve both.
The $AdS$ vector index $a$ carried by the three-fermion piece is what produces the spin
$(\tfrac12,\tfrac12)$ on each side.

\paragraph{RR states with $\langle 4,0\rangle_{(1,1)}$.}
From \eqref{128decomp}, the only $SO(5)$ representation in the ${\bf 128}$ whose tensor
square contains $\langle 4,0\rangle$ is $\langle 1,1\rangle$:
\be
\langle 1,1\rangle\otimes\langle 1,1\rangle\supset\langle 4,0\rangle.
\ee
But the $\langle 1,1\rangle$ component of the ${\bf 128}$ only carries $AdS$ spins
$(\tfrac12,0)$ and $(0,\tfrac12)$. The possible products are:

\smallskip
\begin{center}
\begin{tabular}{|l|l|c|}
\hline
Left spin & Right spin & $(1,1)$? \\
\hline
$(\tfrac12,0)\otimes(\tfrac12,0)=(0,0)\oplus(1,0)$ && No \\
$(\tfrac12,0)\otimes(0,\tfrac12)=(\tfrac12,\tfrac12)$ && No \\
$(0,\tfrac12)\otimes(\tfrac12,0)=(\tfrac12,\tfrac12)$ && No \\
$(0,\tfrac12)\otimes(0,\tfrac12)=(0,0)\oplus(0,1)$ && No \\
\hline
\end{tabular}
\end{center}
\smallskip

None of these reach spin $(1,1)$.

\smallskip
\noindent\textbf{RR total: 0 states.}

\subsection{Higher-harmonic states and the resolution of $b_0$}
\label{sec:b0resolution}

The $\ell=0$ spectrum in \S\ref{sec:spectrum} is incomplete: the $\langle 4,0\rangle_{(1,1)}$
sector at mass level~1 also contains states with \emph{higher $S^5$ harmonics}.
A state with oscillator $SO(5)$ representation $R_{\rm osc}$ and harmonic $\ell$ has total
$SO(5)$ equal to an irreducible component of $R_{\rm osc}\otimes\langle 0,\ell\rangle$.
Since harmonic polynomials $(n\cdot X)^\ell$ with $n\cdot n=0$ are annihilated by the
flat-space Laplacian, they satisfy the same zeroth-order mass-shell condition as $\ell=0$ states.
All harmonics are therefore exactly degenerate at zeroth order.

\paragraph{$\ell=2$ NSNS states with $\langle 4,0\rangle_{(1,1)}$.}
The $\ell=2$ harmonic $Y_{(2)}\in\langle 0,2\rangle$ opens several new channels.
For the total $SO(5)$ to contain $\langle 4,0\rangle$ one needs
$R_{\rm osc}\otimes\langle 0,2\rangle\supset\langle 4,0\rangle$.
For spin $(1,1)$, the available oscillator pairs include:
\begin{itemize}
\item ${\bf 44}\big((1,1),\langle 0,0\rangle\big)\otimes{\bf 44}\big((0,0),\langle 0,2\rangle\big)$:
spin $(1,1)\otimes(0,0)=(1,1)$. Osc.\ $SO(5)$: $\langle 0,2\rangle$.
Then $\langle 0,2\rangle\otimes\langle 0,2\rangle_{\rm harm}\supset\langle 4,0\rangle$. \checkmark
\item ${\bf 44}\big((1,1),\langle 0,0\rangle\big)\otimes{\bf 84}\big((0,0),\langle 2,0\rangle\big)$:
spin $(1,1)$. Osc.\ $SO(5)$: $\langle 2,0\rangle$.
Then $\langle 2,0\rangle\otimes\langle 0,2\rangle\supset\langle 4,0\rangle$. \checkmark
\item ${\bf 84}\big((1,0),\langle 0,1\rangle\big)\otimes{\bf 84}\big((0,1),\langle 0,1\rangle\big)$:
spin $(1,0)\otimes(0,1)=(1,1)$. Osc.\ $SO(5)$:
$\langle 0,1\rangle\otimes\langle 0,1\rangle\supset\langle 2,0\rangle$.
Then $\langle 2,0\rangle\otimes\langle 0,2\rangle\supset\langle 4,0\rangle$. \checkmark
\end{itemize}
Together with mirrors, this gives at least \textbf{6 additional NSNS states at $\ell=2$}.

\paragraph{$\ell=2$ RR states with $\langle 4,0\rangle_{(1,1)}$.}
This is the key difference from the $\ell=0$ analysis.
At $\ell=0$, the only $SO(5)$ representation in the ${\bf 128}$ whose square contains
$\langle 4,0\rangle$ is $\langle 1,1\rangle$, and its $AdS$ spins
$(\tfrac12,0)$, $(0,\tfrac12)$ cannot reach $(1,1)$.

At $\ell=2$, lower oscillator $SO(5)$ suffices because the harmonic makes up the difference.
The $\langle 1,0\rangle$ component of the ${\bf 128}$ carries $AdS$ spins
$(1,\tfrac12)$, $(\tfrac12,1)$, $(\tfrac12,0)$, $(0,\tfrac12)$
(from eq.~\eqref{128decomp}). The products
\be
(1,\tfrac12)\otimes(0,\tfrac12)
\supset(1,1),
\qquad
(\tfrac12,1)\otimes(\tfrac12,0)
\supset(1,1)
\ee
\emph{do} contain spin $(1,1)$.
The oscillator $SO(5)$ product $\langle 1,0\rangle\otimes\langle 1,0\rangle
=\langle 2,0\rangle\oplus\langle 0,1\rangle\oplus\langle 0,0\rangle$
contains $\langle 2,0\rangle$, and
$\langle 2,0\rangle\otimes\langle 0,2\rangle_{\rm harm}\supset\langle 4,0\rangle$.

\smallskip
\noindent\textbf{At least 2 RR states at $\ell=2$ with $\langle 4,0\rangle_{(1,1)}$.}

\paragraph{Corrected spectrum.}
\begin{center}
\begin{tabular}{|l|c|c|}
\hline
& NSNS & RR \\
\hline
$\ell=0$ & 1 & 0 \\
$\ell=2$ & $\geq 6$ & $\geq 2$ \\
\hline
Total & $\geq 7$ & $\geq 2$ \\
\hline
\end{tabular}
\end{center}

\paragraph{Resolution.}
With both NSNS and RR states now present in the $\langle 4,0\rangle_{(1,1)}$ sector
(at $\ell=2$), the off-diagonal first-order perturbation matrix is nontrivial.
The structure is
\be
M=\begin{pmatrix}
0 & V \\ W & 0
\end{pmatrix},
\ee
where $V$ couples NSNS states to the bracket $[\psi^{(1)}_{RR}\otimes\text{RR}]'$
(which is NSNS-valued) and $W$ couples RR states to
$[\psi^{(1)}_{RR}\otimes\text{NSNS}]'$ (which is RR-valued).
Diagonalising $M$ in the full sector
(including $\ell=0$ and $\ell=2$ states simultaneously) produces eigenvalues that should
reproduce $b_0=\D_0-4$ for each physical state.

The physical Konishi state with $[2,0,2]_{(1,1)}$ quantum numbers at $\D_0=6$ is therefore
a specific superposition of the $\ell=0$ NSNS state \eqref{ell0_8484} and the $\ell=2$ NSNS
and RR states listed above, with $b_0=+2$ emerging as an eigenvalue of $M$.

\paragraph{Which harmonic levels contribute?}
A mass-level-1 state with combined oscillator $SO(5)$ representation $R_{\rm osc}$ and
harmonic $\ell$ contributes to total $\langle 4,0\rangle$ iff
$\langle 4,0\rangle\subset R_{\rm osc}\otimes\langle 0,\ell\rangle$.
By reciprocity this is the same as $R_{\rm osc}\subset\langle 4,0\rangle\otimes\langle 0,\ell\rangle$.
The available combined oscillator reps at mass level~1 are finite, so the tower terminates at
a finite $\ell_{\rm max}$.

A key simplification: products of the form $\langle 0,a\rangle\otimes\langle 0,\ell\rangle$
never produce $\langle 4,0\rangle$. This is because these products only contain irreps
$\langle 2k,c\rangle$ with $k\leq\min(a,\ell)$, and reaching $k=2$ (i.e.\ first Dynkin label
$\geq 4$) requires both $a\geq 2$ and $\ell\geq 2$.  For $\langle 0,0\rangle$ and
$\langle 0,1\rangle$ oscillator content this condition is never met.

The surviving combined oscillator reps and their minimum $\ell$ are:

\begin{center}
\begin{tabular}{|l|c|c|l|}
\hline
Combined osc.\ $R_{\rm osc}$ & Min $\ell$ & $\D_0$ & Source \\
\hline
$\langle 4,0\rangle$ & 0 & 6 &
$\langle 0,2\rangle\!\otimes\!\langle 0,2\rangle$,
$\langle 2,0\rangle\!\otimes\!\langle 2,0\rangle$,
$\langle 0,2\rangle\!\otimes\!\langle 2,0\rangle$ \\
$\langle 2,1\rangle$ & 1 & 5 &
$\langle 0,1\rangle\!\otimes\!\langle 2,0\rangle$,
$\langle 2,0\rangle\!\otimes\!\langle 2,0\rangle$, etc. \\
$\langle 2,2\rangle$ & 1 & 5 &
$\langle 0,2\rangle\!\otimes\!\langle 0,2\rangle$ \\
$\langle 2,0\rangle$ & 2 & 4 &
$\langle 0,1\rangle\!\otimes\!\langle 0,1\rangle$,
$\langle 2,0\rangle\!\otimes\!\langle 0,0\rangle$, etc. \\
$\langle 0,2\rangle$ & 2 & 4 &
$\langle 0,2\rangle\!\otimes\!\langle 0,0\rangle$, etc. \\
$\langle 0,4\rangle$ & 2 & 4 &
$\langle 0,2\rangle\!\otimes\!\langle 0,2\rangle$ \\
$\langle 0,0\rangle$ & --- & --- & never reaches $\langle 4,0\rangle$ \\
$\langle 0,1\rangle$ & --- & --- & never reaches $\langle 4,0\rangle$ \\
$\langle 0,3\rangle$ & --- & --- & not available at spin $(0,0)$ or $(1,1)$ \\
\hline
\end{tabular}
\end{center}

The ``$\D_0$'' column indicates which Konishi multiplet level the corresponding
harmonic is naturally associated with: $\D_0=2+\ell+(\text{osc.\ $S^5$ excitation number})$.

The tower thus runs from $\ell=0$ (the $\D_0=6$ states) through $\ell=1$ ($\D_0=5$) to $\ell=2$
($\D_0=4$). Higher harmonics $\ell\geq 3$ do not contribute to total $\langle 4,0\rangle$
from the available mass-level-1 oscillator content, because the oscillator $SO(5)$ reps
that remain ($\langle 0,0\rangle$, $\langle 0,1\rangle$) can never bridge the gap to
$\langle 4,0\rangle$.

\paragraph{Different $\ell$ do not mix at first order.}
The first-order RR background profile is constant:
$F^{(1)}_{\alpha\beta}=\frac{i}{R}(\gamma^{01234})_{\alpha\beta}$,
with no spacetime dependence. The bracket $[\psi^{(1)}\otimes\varphi^{(0)}]'$ acts on
the oscillator content through the OPE but carries the zero-mode profile through unchanged.
A state with harmonic $Y_\ell(X^i)$ maps to a state with the same harmonic.

The first-order perturbation matrix is therefore \emph{block-diagonal in $\ell$}:
states at different harmonic levels do not mix.

\paragraph{Consequence for $b_0$.}
This means the $\ell=0$ block and the $\ell=2$ block are independent at first order.
Each block has its own eigenvalues:
\begin{itemize}
\item $\ell=0$ block: 1 NSNS, 0 RR. The matrix is $1\times 1$ and vanishes by sector.
Eigenvalue: $b_0=0$.
\item $\ell=1$ block: some NSNS and possibly RR states. Eigenvalues: some set of $b_0$ values.
\item $\ell=2$ block: $\geq 5$ NSNS $+$ $\geq 2$ RR. Nontrivial NSNS--RR mixing. Eigenvalues
include the $\D_0=4$ value $b_0=0$, but there may be additional eigenvalues.
\end{itemize}

The $\ell=0$ state, which should be the $\D_0=6$ Konishi descendant with $b_0=+2$, gets
$b_0=0$ from the first-order perturbation theory. The expected constant correction
$b_0=\D_0-4$ is therefore \emph{not} produced by the first-order NSNS--RR mixing mechanism.

\paragraph{Subleading order: $\ell$ mixing from $\psi^{(2)}$.}
The second-order NSNS background has quadratic $X$-dependence:
\be
h^{(2)}_{\mu\nu}\sim -\frac{1}{R^2}(\delta^a_\mu\delta^b_\nu X_aX_b
-\delta^i_\mu\delta^j_\nu X_iX_j).
\ee
The $X_iX_j$ piece multiplies the harmonic of the state and shifts it by $\pm 2$ via the
Clebsch--Gordan decomposition $\langle 0,2\rangle\otimes\langle 0,\ell\rangle
\to\langle 0,\ell\rangle\oplus\langle 0,\ell\pm 2\rangle\oplus\cdots$.
So the bracket $[\psi^{(2)}\otimes\varphi^{(0)}]'$ mixes $\ell$ with $\ell\pm 2$, and the
second-order Hamiltonian $H^{(2)}=[\psi^{(2)}\otimes H^{(0)}]'+\cdots$ acts across harmonic levels.

Since $H^{(1)}=0$, the full spectrum --- both the constant $b_0=\D_0-4$ and the coefficient
$b_1$ --- must emerge from $H^{(2)}$ acting on the combined $\ell=0,1,2$ space.
The $\ell=0$ state ($\D_0=6$) communicates with the $\ell=2$ states ($\D_0=4$) through
the NSNS coupling mediated by $h^{(2)}_{ij}\sim X_iX_j$, so no sector mismatch arises.

Whether the eigenvalue of $H^{(2)}$ in this space is $O(1)$ (giving $b_0$) or
$O(1/R)$ (giving $b_1/R$) depends on the $R$-dependent normalization of the states
and the Hamiltonian, which is what the second-order computation in \texttt{draftKonishi.tex}
is working out.

\section{Mass-level-1 spectrum for the \texorpdfstring{$\D_0=4$}{Delta0=4} candidate}
\label{sec:delta4_spectrum}

We now perform for the $\D_0=4$ candidate $[2,0,2]_{(0,0)}$ the same type of spectrum
survey as for the $\D_0=6$ candidate in \S\ref{sec:spectrum}, with one important
complication: the $\D_0=4$ vertex \eqref{202vertex} carries an $S^5$ harmonic $Y_{(2)}(U)$,
so the total $SO(5)$ quantum number of the state receives contributions from both the
oscillator content and the harmonic profile.

\subsection{Total \texorpdfstring{$SO(5)$}{SO(5)} of the vertex \eqref{202vertex}}

Under $SO(5)$, the separate pieces of the vertex transform as:
\begin{itemize}
\item Profile $Y_{(2)}(U)=(U^+)^2/R^2$: a rank-2 harmonic, carrying
$\langle 0,2\rangle$ with highest weight $2e_1=(2,0)$.
\item Left oscillators $\psi^{V^+}\partial V^+$ (both in $S^5$): a symmetric tensor in $S^5$
indices, sitting in the $\langle 0,2\rangle$ part of $\langle 0,1\rangle\otimes\langle 0,1\rangle$,
with highest weight $2e_2=(0,2)$.
\item Right oscillators $\bar\psi^{V^+}\bar\partial V^+$: same, weight $(0,2)$.
\end{itemize}
The total $SO(5)$ weight of the highest-weight component is
\be
(2,0)+(0,2)+(0,2)=(2,4).
\ee
This is \emph{not} a dominant weight of $SO(5)$ (in Dynkin labels it
would require $b=2-4<0$). The resolution is that $\psi^{V^+}\partial V^+$
is not purely in $\langle 0,2\rangle$: its $SO(5)$ decomposition
$\langle 0,1\rangle\otimes\langle 0,1\rangle=\langle 0,0\rangle\oplus\langle 0,2\rangle\oplus\langle 2,0\rangle$
also includes a trace $\langle 0,0\rangle$ and an antisymmetric part $\langle 2,0\rangle$.
The physical vertex is the BRST-closed combination, which involves all three pieces
(plus the three-fermion and $\partial\psi$ completions required by BRST closure).

\paragraph{Identifying the total $SO(5)$ irrep.}
The $[2,0,2]$ representation of $SO(6)$ branches as $\langle 4,0\rangle\oplus\langle 2,1\rangle
\oplus\langle 0,2\rangle$ under $SO(5)$. The vertex \eqref{202vertex} is the
$SO(5)$ highest-weight component of the Konishi state, and the top branch of
$[2,0,2]$ is $\langle 4,0\rangle$ (highest weight $(2,2)$). The dominant weight $(2,2)$
is reachable from the profile weight $(2,0)$ combined with the $\langle 2,0\rangle$
antisymmetric part of each oscillator pair (weight $(1,1)$ on left, $(1,1)$ on right is
too much; more precisely the $\langle 4,0\rangle$ is assembled from
$\langle 0,2\rangle_{\rm profile}\otimes\langle 2,0\rangle_{\rm osc,L}
\otimes\langle 2,0\rangle_{\rm osc,R}$, which does contain $\langle 4,0\rangle$
via the chain $\langle 2,0\rangle\otimes\langle 2,0\rangle\supset\langle 4,0\rangle$
and then $\langle 4,0\rangle\otimes\langle 0,2\rangle\supset\langle 4,0\rangle$).

The total $SO(5)$ representation of the $[2,0,2]_{(0,0)}$ Konishi state
at $\D_0=4$ is thus $\langle 4,0\rangle$, with boundary spin $(0,0)$.

\subsection{NSNS states in the \texorpdfstring{$\langle 4,0\rangle_{(0,0)}$}{<4,0>(0,0)} sector}

A mass-level-1 NSNS state with total $SO(5)=\langle 4,0\rangle$ and spin $(0,0)$
is built from an oscillator content (an $SO(9)$ representation on each side) together with
an $S^5$ harmonic of level~$\ell$. The total $SO(5)$ is the irreducible component
$\langle 4,0\rangle$ inside (oscillator $SO(5)$)$\,\otimes\,\langle 0,\ell\rangle$.

From the decompositions \eqref{44decomp} and \eqref{84decomp}, the spin-$(0,0)$
oscillator $SO(5)$ representations are
\be
{\bf 44}:\;\langle 0,0\rangle,\;\langle 0,2\rangle;
\qquad
{\bf 84}:\;\langle 2,0\rangle.
\ee
For the closed string we take left$\,\otimes\,$right, giving $3\times 3=9$ possible
oscillator pairs. Denoting the oscillator product as $R_{\rm osc}$, we need
$R_{\rm osc}\otimes\langle 0,\ell\rangle\supset\langle 4,0\rangle$.

\paragraph{States with $\ell=0$ (constant profile on $S^5$).}
The total $SO(5)$ equals the oscillator $SO(5)$. Using the results from the
$\D_0=6$ analysis (\S\ref{sec:spectrum}):
\begin{center}
\begin{tabular}{|l|l|c|}
\hline
Sector & Osc.\ $SO(5)$ product & $\langle 4,0\rangle$? \\
\hline
${\bf 44}\otimes{\bf 44}$ & $\langle 0,2\rangle\otimes\langle 0,2\rangle\supset\langle 4,0\rangle$ & Yes \\
${\bf 44}\otimes{\bf 84}$ & $\langle 0,2\rangle\otimes\langle 2,0\rangle\supset\langle 4,0\rangle$ & Yes \\
${\bf 84}\otimes{\bf 44}$ & $\langle 2,0\rangle\otimes\langle 0,2\rangle\supset\langle 4,0\rangle$ & Yes \\
${\bf 84}\otimes{\bf 84}$ & $\langle 2,0\rangle\otimes\langle 2,0\rangle\supset\langle 4,0\rangle$ & Yes \\
\hline
\end{tabular}
\end{center}
\textbf{4 NSNS states at $\ell=0$.}

\paragraph{States with $\ell=2$.}
Now we need $R_{\rm osc}\otimes\langle 0,2\rangle\supset\langle 4,0\rangle$.
The additional oscillator products that work (those that did not already give
$\langle 4,0\rangle$ at $\ell=0$):
\begin{itemize}
\item $\langle 0,0\rangle\otimes\langle 0,2\rangle=\langle 0,2\rangle$, then
$\langle 0,2\rangle\otimes\langle 0,2\rangle\supset\langle 4,0\rangle$. Sources:
${\bf 44}(\langle 0,0\rangle)\otimes{\bf 44}(\langle 0,2\rangle)$, its mirror,
${\bf 44}(\langle 0,0\rangle)\otimes{\bf 84}(\langle 2,0\rangle)$, its mirror.
\item $\langle 0,2\rangle\otimes\langle 0,0\rangle$: same as above.
\item $\langle 2,0\rangle\otimes\langle 0,0\rangle=\langle 2,0\rangle$, then
$\langle 2,0\rangle\otimes\langle 0,2\rangle\supset\langle 4,0\rangle$.
\item The 4 products from $\ell=0$ also produce $\langle 4,0\rangle$ at $\ell=2$
through subleading components.
\end{itemize}
Counting carefully (each of $\langle 0,0\rangle$, $\langle 0,2\rangle$, $\langle 2,0\rangle$
on each side, combined in pairs): at least \textbf{5 additional NSNS states at $\ell=2$}
(the exact count depends on which tensor-product components are independent after
BRST projection).

Higher harmonics $\ell=4,6,\ldots$ contribute further states.

\subsection{RR states in the \texorpdfstring{$\langle 4,0\rangle_{(0,0)}$}{<4,0>(0,0)} sector}

This is the critical difference from the $\D_0=6$ analysis.

From \eqref{128decomp}, the $\langle 1,1\rangle$ component of the ${\bf 128}$ carries
$AdS$ spins $(\tfrac12,0)$ and $(0,\tfrac12)$. The product
\be
(\tfrac12,0)\otimes(\tfrac12,0)=(0,0)\oplus(1,0),
\qquad
(0,\tfrac12)\otimes(0,\tfrac12)=(0,0)\oplus(0,1)
\ee
\emph{does} contain spin $(0,0)$. (Contrast the $\D_0=6$ case: spin $(1,1)$ was
unreachable from these half-integer spins.)

The oscillator $SO(5)$ product $\langle 1,1\rangle\otimes\langle 1,1\rangle$ contains
$\langle 4,0\rangle$ (as verified in \S\ref{sec:spectrum}). Combined with $\ell=0$ harmonic,
this gives total $\langle 4,0\rangle$.

\smallskip\noindent
\textbf{RR states at $\ell=0$:}
\begin{center}
\begin{tabular}{|l|l|c|}
\hline
Left ${\bf 128}$ & Right ${\bf 128}$ & spin $(0,0)$? \\
\hline
$(\tfrac12,0),\langle 1,1\rangle$ & $(\tfrac12,0),\langle 1,1\rangle$ & Yes \\
$(0,\tfrac12),\langle 1,1\rangle$ & $(0,\tfrac12),\langle 1,1\rangle$ & Yes \\
$(\tfrac12,0),\langle 1,1\rangle$ & $(0,\tfrac12),\langle 1,1\rangle$ & No (spin $(\tfrac12,\tfrac12)$)\\
$(0,\tfrac12),\langle 1,1\rangle$ & $(\tfrac12,0),\langle 1,1\rangle$ & No \\
\hline
\end{tabular}
\end{center}
\textbf{2 RR states at $\ell=0$.}

Higher harmonics contribute further RR states (e.g.\ $\langle 1,0\rangle\otimes\langle 1,0\rangle$
oscillator content combined with $\ell\geq 4$ harmonic).

\subsection{Explicit vertex operators}
\label{sec:explicit_vertices}

We list the highest-weight forms of all states in the $\langle 4,0\rangle_{(0,0)}$
sector at mass level~1 with $\ell\leq 2$, using the complex coordinates
$U^\pm=X^5\pm iX^6$, $V^\pm=X^7\pm iX^8$, and $X^9$ for the $S^5$ directions.
Each vertex is written in the $(-1,-1)$ picture (NSNS) or $(-\tfrac12,-\tfrac12)$
picture (RR); BRST completions (from $\partial\psi$, gamma-trace subtractions, etc.)
are suppressed but required for the full physical state.

\subsubsection*{$\ell=0$ NSNS (4 states)}

These have constant $S^5$ profile. The $\langle 4,0\rangle$ at highest weight $(2,2)$
is assembled entirely from oscillators; the two worldsheet sectors must use
\emph{different} $S^5$ planes.

\paragraph{${\bf 44}\otimes{\bf 44}$:}
Left oscillator in $\langle 0,2\rangle$ at weight $(2,0)$, right at $(0,2)$:
\be\label{ell0_4444}
\varphi^{(\ell=0)}_{44\times 44}
= a(X^a)\, c\bar c\,
e^{-\phi}\psi^{U^+}\partial U^+\,
e^{-\bar\phi}\bar\psi^{V^+}\bar\partial V^+
+\cdots
\ee
Note the left oscillator is in the $U$-plane and the right in the $V$-plane.

\paragraph{${\bf 44}\otimes{\bf 84}$:}
Left $\langle 0,2\rangle$ at weight $(1,1)$, right $\langle 2,0\rangle$ at weight $(1,1)$:
\be\label{ell0_4484}
\varphi^{(\ell=0)}_{44\times 84}
= a(X^a)\, c\bar c\,
e^{-\phi}\big(\psi^{U^+}\partial V^++\psi^{V^+}\partial U^+\big)\,
e^{-\bar\phi}\bar\psi^{U^+}\bar\psi^{V^+}\bar\psi^9
+\cdots
\ee
The left side is the symmetric combination in the $U$--$V$ plane;
the right side is the three-fermion $\Lambda^3(S^5)$ vertex.

\paragraph{${\bf 84}\otimes{\bf 44}$:}
Mirror of \eqref{ell0_4484}:
\be\label{ell0_8444}
\varphi^{(\ell=0)}_{84\times 44}
= a(X^a)\, c\bar c\,
e^{-\phi}\psi^{U^+}\psi^{V^+}\psi^9\,
e^{-\bar\phi}\big(\bar\psi^{U^+}\bar\partial V^++\bar\psi^{V^+}\bar\partial U^+\big)
+\cdots
\ee

\paragraph{${\bf 84}\otimes{\bf 84}$:}
Both sides three-fermion, $\langle 2,0\rangle\otimes\langle 2,0\rangle$ at weight $(1,1)+(1,1)=(2,2)$:
\be\label{ell0_8484}
\varphi^{(\ell=0)}_{84\times 84}
= a(X^a)\, c\bar c\,
e^{-\phi}\psi^{U^+}\psi^{V^+}\psi^9\,
e^{-\bar\phi}\bar\psi^{U^+}\bar\psi^{V^+}\bar\psi^9
+\cdots
\ee

\subsubsection*{$\ell=0$ RR (2 states)}

Both use the $\langle 1,1\rangle$ component of the ${\bf 128}$ (gamma-traceless
combination of $S^\alpha\partial X^i$ and $S^i_\alpha$) on each side, at the $\langle 4,0\rangle$
highest weight. The $\langle 1,1\rangle$ on the left carries the $(\tfrac12,\tfrac32)$
weight of $\langle 1,1\rangle$ (spin field in $U^+$ with excitation in $V^+$),
and the right carries $(\tfrac32,\tfrac12)$ (spin field in $V^+$ with excitation in $U^+$),
or vice versa.

\paragraph{RR, chirality $(+,+)$:}
Both sides with $AdS$ spin $(\tfrac12,0)$:
\be\label{ell0_RR1}
\varphi^{(\ell=0)}_{RR,+}
= g(X^a)\, c\bar c\,
e^{-\frac{1}{2}\phi}\big(S^\alpha\partial V^+
-\text{$\gamma$-tr.}\big)_{\langle 1,1\rangle}\,
e^{-\frac{1}{2}\bar\phi}\big(\bar S^\beta\bar\partial U^+
-\text{$\gamma$-tr.}\big)_{\langle 1,1\rangle}
+\cdots
\ee
where $\alpha,\beta$ are projected onto the $(\tfrac12,0)$ $AdS$ chirality, and the
bispinor profile $g(X^a)$ contracts the spinor indices appropriately.

\paragraph{RR, chirality $(-,-)$:}
Both sides with $AdS$ spin $(0,\tfrac12)$:
\be\label{ell0_RR2}
\varphi^{(\ell=0)}_{RR,-}
= \tilde g(X^a)\, c\bar c\,
e^{-\frac{1}{2}\phi}\big(S^\alpha\partial V^+
-\text{$\gamma$-tr.}\big)_{\langle 1,1\rangle}\,
e^{-\frac{1}{2}\bar\phi}\big(\bar S^\beta\bar\partial U^+
-\text{$\gamma$-tr.}\big)_{\langle 1,1\rangle}
+\cdots
\ee
with $\alpha,\beta$ in the $(0,\tfrac12)$ $AdS$ chirality.

\subsubsection*{$\ell=2$ NSNS}

These carry the harmonic $Y_{(2)}(U)=(U^+)^2/R^2\in\langle 0,2\rangle$.
The $\langle 4,0\rangle$ is assembled from the harmonic's $\langle 0,2\rangle$
combined with the oscillator content.

\paragraph{${\bf 44}\otimes{\bf 44}$, both $\langle 0,2\rangle$ (the Konishi seed):}
This is the original vertex \eqref{202vertex}:
\be\label{ell2_4444a}
\varphi^{(\ell=2)}_{44\times 44,\, a}
= Y_{(2)}(U)\, c\bar c\,
e^{-\phi}\psi^{V^+}\partial V^+\,
e^{-\bar\phi}\bar\psi^{V^+}\bar\partial V^+
+\cdots
\ee
Both oscillators are in the $V$-plane; the $U$-plane charge comes entirely from the harmonic.

\paragraph{${\bf 44}\otimes{\bf 44}$, $\langle 0,2\rangle\otimes\langle 0,0\rangle$:}
Right oscillator is the scalar piece of the ${\bf 44}$ (the traceless combination
$\eta^{ij}\psi^i\partial X^j$ with trace subtracted against the $AdS$ piece):
\be\label{ell2_4444b}
\varphi^{(\ell=2)}_{44\times 44,\, b}
= Y_{(2)}(U)\, c\bar c\,
e^{-\phi}\psi^{V^+}\partial V^+\,
e^{-\bar\phi}\bar\psi^a\bar\partial X_a\big|_{\rm traceless}
+\cdots
\ee

\paragraph{${\bf 44}\otimes{\bf 44}$, $\langle 0,0\rangle\otimes\langle 0,2\rangle$:}
Mirror of \eqref{ell2_4444b}.

\paragraph{${\bf 44}\otimes{\bf 84}$, $\langle 0,0\rangle\otimes\langle 2,0\rangle$:}
\be\label{ell2_4484}
\varphi^{(\ell=2)}_{44\times 84}
= Y_{(2)}(U)\, c\bar c\,
e^{-\phi}\psi^a\partial X_a\big|_{\rm traceless}\,
e^{-\bar\phi}\bar\psi^{U^+}\bar\psi^{V^+}\bar\psi^9
+\cdots
\ee
and its mirror ${\bf 84}\otimes{\bf 44}$.

\paragraph{${\bf 84}\otimes{\bf 84}$, $\langle 2,0\rangle\otimes\langle 2,0\rangle$ with $\ell=2$:}
\be\label{ell2_8484}
\varphi^{(\ell=2)}_{84\times 84}
= Y_{(2)}(U)\, c\bar c\,
e^{-\phi}\psi^{U^+}\psi^{V^+}\psi^9\,
e^{-\bar\phi}\bar\psi^{U^+}\bar\psi^{V^+}\bar\psi^9
+\cdots
\ee
This is the $\ell=0$ state \eqref{ell0_8484} dressed with the $\ell=2$ harmonic;
the $\langle 4,0\rangle$ projection differs from the $\ell=0$ case because
the harmonic shifts the $SO(5)$ content.

\paragraph{Remark.}
The $\ell=2$ Konishi seed \eqref{ell2_4444a} is not a pure $\langle 4,0\rangle$ eigenstate.
The component written has total $SO(5)$ weight
$(2,0)+(0,2)+(0,2)=(2,4)$, which exceeds the highest weight $(2,2)$ of $\langle 4,0\rangle$.
The full $\langle 4,0\rangle$ projection requires a specific linear combination of oscillator
components --- including the $\langle 2,0\rangle$ (antisymmetric) and $\langle 0,0\rangle$
(trace) parts of $\psi^i\partial X^j$ alongside the symmetric $\langle 0,2\rangle$ --- that
is determined by BRST closure. The highest-weight forms above should therefore be understood
as schematic leading components, not as complete physical vertices.

\subsection{Summary for the \texorpdfstring{$\D_0=4$}{Delta0=4} sector}

The $\langle 4,0\rangle_{(0,0)}$ sector at mass level~1 is much richer than the
$\langle 4,0\rangle_{(1,1)}$ sector:
\begin{center}
\begin{tabular}{|l|c|c|}
\hline
& $\langle 4,0\rangle_{(1,1)}$ ($\D_0=6$) & $\langle 4,0\rangle_{(0,0)}$ ($\D_0=4$) \\
\hline
NSNS ($\ell=0$) & 1 & 4 \\
NSNS ($\ell=2$) & --- & $\geq 5$ \\
RR ($\ell=0$) & 0 & 2 \\
\hline
Total (low $\ell$) & \textbf{1} & $\geq\textbf{11}$ \\
\hline
\end{tabular}
\end{center}

There are three important consequences.
\begin{enumerate}
\item \textbf{The zeroth-order state is pure NSNS.}
In flat space the NSNS and RR sectors decouple.  The Konishi state at $\D_0=4$ starts life
as a single NSNS state; RR admixtures are generated perturbatively
by the first-order RR background.  Since $b_0=0$ there is no first-order splitting that
would require a mixed zeroth-order ansatz.

\item \textbf{The sector is degenerate.}
Multiple NSNS and RR states share the quantum numbers $\langle 4,0\rangle_{(0,0)}$ at
mass level~1. In the absence of a first-order splitting (since $H^{(1)}=0$), these
states remain degenerate through first order. Their degeneracy is lifted at
second order by the NSNS background $\psi^{(2)}$ and by the iterated first-order
RR correction.

\item \textbf{The second-order computation is a degenerate perturbation theory problem.}
To extract $b_1=2$ one must diagonalise the second-order perturbation matrix in the full
$\langle 4,0\rangle_{(0,0)}$ sector. The simple one-vertex ansatz \eqref{202vertex}
is a good \emph{starting point} (it captures the leading oscillator and harmonic structure),
but the physical Konishi eigenstate at second order will generically be a superposition of
the states enumerated above.
\end{enumerate}

\appendix

\section{Explicit \texorpdfstring{$SO(6)\to SO(5)$}{SO(6)->SO(5)} Branching Data}
\label{app:branching}

For the explicit Table-1 scan I use the following $SO(6)\to SO(5)$ branchings, written in the
$C_2\simeq B_2$ convention where $\langle 1,0\rangle={\bf 4}$ and $\langle 0,1\rangle={\bf 5}$.
The pairs related by conjugation in $SU(4)\simeq SO(6)$ branch identically to $SO(5)$.

\begin{align*}
[0,0,0] &\to \langle 0,0\rangle, \\
[0,0,1]=[1,0,0] &\to \langle 1,0\rangle, \\
[0,1,0] &\to \langle 0,1\rangle\oplus \langle 0,0\rangle, \\
[0,0,2]=[2,0,0] &\to \langle 2,0\rangle, \\
[1,0,1] &\to \langle 2,0\rangle\oplus \langle 0,1\rangle, \\
[0,1,1]=[1,1,0] &\to \langle 1,1\rangle\oplus \langle 1,0\rangle, \\
[0,0,3]=[3,0,0] &\to \langle 3,0\rangle, \\
[1,0,2]=[2,0,1] &\to \langle 3,0\rangle\oplus \langle 1,1\rangle, \\
[0,2,0] &\to \langle 0,2\rangle\oplus \langle 0,1\rangle\oplus \langle 0,0\rangle, \\
[0,1,2]=[2,1,0] &\to \langle 2,1\rangle\oplus \langle 2,0\rangle, \\
[1,1,1] &\to \langle 2,1\rangle\oplus \langle 2,0\rangle\oplus \langle 0,2\rangle\oplus \langle 0,1\rangle, \\
[0,2,1]=[1,2,0] &\to \langle 1,2\rangle\oplus \langle 1,1\rangle\oplus \langle 1,0\rangle, \\
[2,1,1]=[1,1,2] &\to \langle 3,1\rangle\oplus \langle 3,0\rangle\oplus \langle 1,2\rangle\oplus \langle 1,1\rangle, \\
[0,3,0] &\to \langle 0,3\rangle\oplus \langle 0,2\rangle\oplus \langle 0,1\rangle\oplus \langle 0,0\rangle, \\
[1,0,3]=[3,0,1] &\to \langle 4,0\rangle\oplus \langle 2,1\rangle, \\
[2,0,2] &\to \langle 4,0\rangle\oplus \langle 2,1\rangle\oplus \langle 0,2\rangle, \\
[0,2,2]=[2,2,0] &\to \langle 2,2\rangle\oplus \langle 2,1\rangle\oplus \langle 2,0\rangle, \\
[1,2,1] &\to \langle 2,2\rangle\oplus \langle 2,1\rangle\oplus \langle 2,0\rangle\oplus \langle 0,3\rangle\oplus \langle 0,2\rangle\oplus \langle 0,1\rangle, \\
[1,3,0]=[0,3,1] &\to \langle 1,3\rangle\oplus \langle 1,2\rangle\oplus \langle 1,1\rangle\oplus \langle 1,0\rangle, \\
[3,1,0]=[0,1,3] &\to \langle 3,1\rangle\oplus \langle 3,0\rangle, \\
[0,0,4]=[4,0,0] &\to \langle 4,0\rangle, \\
[0,4,0] &\to \langle 0,4\rangle\oplus \langle 0,3\rangle\oplus \langle 0,2\rangle\oplus \langle 0,1\rangle\oplus \langle 0,0\rangle.
\end{align*}

\section{Table 1 of \texttt{1102.1209}}
\label{app:table1}

For convenience I reproduce Table 1 of \texttt{1102.1209} here, with the original
shorthand written out explicitly in terms of $\Delta_0$.

\footnotesize
\begin{longtable}{|c|p{0.77\textwidth}|}
\hline
$\D_0$ & States \\
\hline
\endfirsthead
\hline
$\D_0$ & States \\
\hline
\endhead
$2$ & $[0,0,0]_{(0,0)}$ \\
\hline
$2+\frac{1}{2}$ & $[0,0,1]_{(0,\frac{1}{2})}+[1,0,0]_{(\frac{1}{2},0)}$ \\
\hline
$2+1$ & $[0,0,0]_{(\frac{1}{2},\frac{1}{2})}+[0,0,2]_{(0,0)}+[0,1,0]_{(0,1)+(1,0)}+[1,0,1]_{(\frac{1}{2},\frac{1}{2})}+[2,0,0]_{(0,0)}$ \\
\hline
$2+\frac{3}{2}$ & $[0,0,1]_{(\frac{1}{2},0)+(\frac{1}{2},1)+(\frac{3}{2},0)}+[0,1,1]_{(0,\frac{1}{2})+(1,\frac{1}{2})}+[1,0,0]_{(0,\frac{1}{2})+(0,\frac{3}{2})+(1,\frac{1}{2})}+[1,0,2]_{(\frac{1}{2},0)}$ \\
& $+[1,1,0]_{(\frac{1}{2},0)+(\frac{1}{2},1)}+[2,0,1]_{(0,\frac{1}{2})}$ \\
\hline
$2+2$ & $[0,0,0]_{(0,0)+(0,2)+(1,1)+(2,0)}+[0,0,2]_{(\frac{1}{2},\frac{1}{2})+(\frac{3}{2},\frac{1}{2})}+[0,1,0]_{2(\frac{1}{2},\frac{1}{2})+(\frac{1}{2},\frac{3}{2})+(\frac{3}{2},\frac{1}{2})}+[2,0,2]_{(0,0)}+[2,1,0]_{(0,1)}$ \\
& $+[0,1,2]_{(1,0)}+[0,2,0]_{2(0,0)+(1,1)}+[1,0,1]_{(0,0)+2(0,1)+2(1,0)+(1,1)}+[1,1,1]_{2(\frac{1}{2},\frac{1}{2})}+[2,0,0]_{(\frac{1}{2},\frac{1}{2})+(\frac{1}{2},\frac{3}{2})}$ \\
\hline
$2+\frac{5}{2}$ & $[0,0,1]_{(0,\frac{1}{2})+(0,\frac{3}{2})+2(1,\frac{1}{2})+(1,\frac{3}{2})+(2,\frac{1}{2})}+[0,0,3]_{(\frac{3}{2},0)}+[0,1,1]_{3(\frac{1}{2},0)+2(\frac{1}{2},1)+(\frac{3}{2},0)+(\frac{3}{2},1)}+[0,2,1]_{(0,\frac{1}{2})+(1,\frac{1}{2})}$ \\
& $+[1,0,0]_{(\frac{1}{2},0)+2(\frac{1}{2},1)+(\frac{1}{2},2)+(\frac{3}{2},0)+(\frac{3}{2},1)}+[1,0,2]_{(0,\frac{1}{2})+2(1,\frac{1}{2})}+[1,1,0]_{3(0,\frac{1}{2})+(0,\frac{3}{2})+2(1,\frac{1}{2})+(1,\frac{3}{2})}$ \\
& $+[1,1,2]_{(\frac{1}{2},0)}+[1,2,0]_{(\frac{1}{2},0)+(\frac{1}{2},1)}+[2,0,1]_{(\frac{1}{2},0)+2(\frac{1}{2},1)}+[2,1,1]_{(0,\frac{1}{2})}+[3,0,0]_{(0,\frac{3}{2})}$ \\
\hline
$2+3$ & $[0,0,0]_{(\frac{1}{2},\frac{1}{2})+(\frac{1}{2},\frac{3}{2})+(\frac{3}{2},\frac{1}{2})+(\frac{3}{2},\frac{3}{2})}+[0,0,2]_{2(0,0)+(1,0)+2(1,1)+(2,0)}+[0,1,0]_{3(0,1)+3(1,0)+2(1,1)+(1,2)+(2,1)}$ \\
& $+[0,1,2]_{2(\frac{1}{2},\frac{1}{2})+(\frac{3}{2},\frac{1}{2})}+[0,2,0]_{3(\frac{1}{2},\frac{1}{2})+(\frac{1}{2},\frac{3}{2})+(\frac{3}{2},\frac{1}{2})}+[0,2,2]_{(0,0)}+[0,3,0]_{(0,1)+(1,0)}$ \\
& $+[1,0,3]_{(1,0)}+[1,1,1]_{2(0,0)+2(0,1)+2(1,0)+2(1,1)}+[1,2,1]_{(\frac{1}{2},\frac{1}{2})}+[2,0,0]_{2(0,0)+(0,1)+(0,2)+2(1,1)}$ \\
& $+[2,0,2]_{(\frac{1}{2},\frac{1}{2})}+[2,1,0]_{2(\frac{1}{2},\frac{1}{2})+(\frac{1}{2},\frac{3}{2})}+[2,2,0]_{(0,0)}+[3,0,1]_{(0,1)}+[1,0,1]_{4(\frac{1}{2},\frac{1}{2})+2(\frac{1}{2},\frac{3}{2})+2(\frac{3}{2},\frac{1}{2})+(\frac{3}{2},\frac{3}{2})}$ \\
\hline
$2+\frac{7}{2}$ & $[0,0,1]_{2(\frac{1}{2},0)+3(\frac{1}{2},1)+(\frac{3}{2},0)+2(\frac{3}{2},1)+(\frac{3}{2},2)}+[0,0,3]_{(0,\frac{1}{2})+(1,\frac{1}{2})}+[0,1,1]_{3(0,\frac{1}{2})+(0,\frac{3}{2})+4(1,\frac{1}{2})+2(1,\frac{3}{2})+(2,\frac{1}{2})}$ \\
& $+[0,1,3]_{(\frac{1}{2},0)}+[0,2,1]_{2(\frac{1}{2},0)+2(\frac{1}{2},1)+(\frac{3}{2},0)}+[0,3,1]_{(0,\frac{1}{2})}+[1,0,0]_{2(0,\frac{1}{2})+(0,\frac{3}{2})+3(1,\frac{1}{2})+2(1,\frac{3}{2})+(2,\frac{3}{2})}$ \\
& $+[1,0,2]_{2(\frac{1}{2},0)+2(\frac{1}{2},1)+(\frac{3}{2},0)+(\frac{3}{2},1)}+[1,1,0]_{3(\frac{1}{2},0)+4(\frac{1}{2},1)+(\frac{1}{2},2)+(\frac{3}{2},0)+2(\frac{3}{2},1)}+[1,1,2]_{(0,\frac{1}{2})+(1,\frac{1}{2})}$ \\
& $+[1,2,0]_{2(0,\frac{1}{2})+(0,\frac{3}{2})+2(1,\frac{1}{2})}+[1,3,0]_{(\frac{1}{2},0)}+[2,0,1]_{2(0,\frac{1}{2})+(0,\frac{3}{2})+2(1,\frac{1}{2})+(1,\frac{3}{2})}+[2,1,1]_{(\frac{1}{2},0)+(\frac{1}{2},1)}$ \\
& $+[3,0,0]_{(\frac{1}{2},0)+(\frac{1}{2},1)}+[3,1,0]_{(0,\frac{1}{2})}$ \\
\hline
$2+4$ & $[0,0,0]_{3(0,0)+3(1,1)+(2,2)}+[0,0,2]_{3(\frac{1}{2},\frac{1}{2})+(\frac{1}{2},\frac{3}{2})+(\frac{3}{2},\frac{1}{2})+(\frac{3}{2},\frac{3}{2})}+[0,1,0]_{4(\frac{1}{2},\frac{1}{2})+2(\frac{1}{2},\frac{3}{2})+2(\frac{3}{2},\frac{1}{2})+2(\frac{3}{2},\frac{3}{2})}$ \\
& $+[0,1,2]_{(0,0)+2(0,1)+2(1,0)+(1,1)}+[0,2,0]_{3(0,0)+(0,1)+(0,2)+(1,0)+3(1,1)+(2,0)}+[0,2,2]_{(\frac{1}{2},\frac{1}{2})}$ \\
& $+[0,3,0]_{2(\frac{1}{2},\frac{1}{2})}+[0,4,0]_{(0,0)}+[1,0,1]_{(0,0)+3(0,1)+3(1,0)+4(1,1)+(1,2)+(2,1)}+[1,0,3]_{(\frac{1}{2},\frac{1}{2})}+[0,0,4]_{(0,0)}$ \\
& $+[1,1,1]_{4(\frac{1}{2},\frac{1}{2})+2(\frac{1}{2},\frac{3}{2})+2(\frac{3}{2},\frac{1}{2})}+[1,2,1]_{(0,0)+(0,1)+(1,0)}+[2,0,0]_{3(\frac{1}{2},\frac{1}{2})+(\frac{1}{2},\frac{3}{2})+(\frac{3}{2},\frac{1}{2})+(\frac{3}{2},\frac{3}{2})}$ \\
& $+[2,0,2]_{(0,0)+(1,1)}+[2,1,0]_{(0,0)+2(0,1)+2(1,0)+(1,1)}+[2,2,0]_{(\frac{1}{2},\frac{1}{2})}+[3,0,1]_{(\frac{1}{2},\frac{1}{2})}+[4,0,0]_{(0,0)}$ \\
\hline
$2+\frac{9}{2}$ & $[0,0,1]_{2(0,\frac{1}{2})+(0,\frac{3}{2})+3(1,\frac{1}{2})+2(1,\frac{3}{2})+(2,\frac{3}{2})}+[0,0,3]_{(\frac{1}{2},0)+(\frac{1}{2},1)}+[0,1,1]_{3(\frac{1}{2},0)+4(\frac{1}{2},1)+(\frac{1}{2},2)+(\frac{3}{2},0)+2(\frac{3}{2},1)}$ \\
& $+[0,1,3]_{(0,\frac{1}{2})}+[0,2,1]_{2(0,\frac{1}{2})+(0,\frac{3}{2})+2(1,\frac{1}{2})}+[0,3,1]_{(\frac{1}{2},0)}+[1,0,0]_{2(\frac{1}{2},0)+3(\frac{1}{2},1)+(\frac{3}{2},0)+2(\frac{3}{2},1)+(\frac{3}{2},2)}$ \\
& $+[1,0,2]_{2(0,\frac{1}{2})+(0,\frac{3}{2})+2(1,\frac{1}{2})+(1,\frac{3}{2})}+[1,1,0]_{3(0,\frac{1}{2})+(0,\frac{3}{2})+4(1,\frac{1}{2})+2(1,\frac{3}{2})+(2,\frac{1}{2})}+[1,1,2]_{(\frac{1}{2},0)+(\frac{1}{2},1)}$ \\
& $+[1,2,0]_{2(\frac{1}{2},0)+2(\frac{1}{2},1)+(\frac{3}{2},0)}+[1,3,0]_{(0,\frac{1}{2})}+[2,0,1]_{2(\frac{1}{2},0)+2(\frac{1}{2},1)+(\frac{3}{2},0)+(\frac{3}{2},1)}+[2,1,1]_{(0,\frac{1}{2})+(1,\frac{1}{2})}$ \\
& $+[3,0,0]_{(0,\frac{1}{2})+(1,\frac{1}{2})}+[3,1,0]_{(\frac{1}{2},0)}$ \\
\hline
$2+5$ & $[0,0,0]_{(\frac{1}{2},\frac{1}{2})+(\frac{1}{2},\frac{3}{2})+(\frac{3}{2},\frac{1}{2})+(\frac{3}{2},\frac{3}{2})}+[0,0,2]_{2(0,0)+(0,1)+(0,2)+2(1,1)}+[0,1,0]_{3(0,1)+3(1,0)+2(1,1)+(1,2)+(2,1)}$ \\
& $+[0,1,2]_{2(\frac{1}{2},\frac{1}{2})+(\frac{1}{2},\frac{3}{2})}+[0,2,0]_{3(\frac{1}{2},\frac{1}{2})+(\frac{1}{2},\frac{3}{2})+(\frac{3}{2},\frac{1}{2})}+[0,2,2]_{(0,0)}+[0,3,0]_{(0,1)+(1,0)}$ \\
& $+[1,0,3]_{(0,1)}+[1,1,1]_{2(0,0)+2(0,1)+2(1,0)+2(1,1)}+[1,2,1]_{(\frac{1}{2},\frac{1}{2})}+[2,0,0]_{2(0,0)+(1,0)+2(1,1)+(2,0)}$ \\
& $+[2,0,2]_{(\frac{1}{2},\frac{1}{2})}+[2,1,0]_{2(\frac{1}{2},\frac{1}{2})+(\frac{3}{2},\frac{1}{2})}+[2,2,0]_{(0,0)}+[3,0,1]_{(1,0)}+[1,0,1]_{4(\frac{1}{2},\frac{1}{2})+2(\frac{1}{2},\frac{3}{2})+2(\frac{3}{2},\frac{1}{2})+(\frac{3}{2},\frac{3}{2})}$ \\
\hline
$2+\frac{11}{2}$ & $[0,0,1]_{(\frac{1}{2},0)+2(\frac{1}{2},1)+(\frac{1}{2},2)+(\frac{3}{2},0)+(\frac{3}{2},1)}+[0,0,3]_{(0,\frac{3}{2})}+[0,1,1]_{3(0,\frac{1}{2})+(0,\frac{3}{2})+2(1,\frac{1}{2})+(1,\frac{3}{2})}+[0,2,1]_{(\frac{1}{2},0)+(\frac{1}{2},1)}$ \\
& $+[1,0,0]_{(0,\frac{1}{2})+(0,\frac{3}{2})+2(1,\frac{1}{2})+(1,\frac{3}{2})+(2,\frac{1}{2})}+[1,0,2]_{(\frac{1}{2},0)+2(\frac{1}{2},1)}+[1,1,0]_{3(\frac{1}{2},0)+2(\frac{1}{2},1)+(\frac{3}{2},0)+(\frac{3}{2},1)}$ \\
& $+[1,1,2]_{(0,\frac{1}{2})}+[1,2,0]_{(0,\frac{1}{2})+(1,\frac{1}{2})}+[2,0,1]_{(0,\frac{1}{2})+2(1,\frac{1}{2})}+[2,1,1]_{(\frac{1}{2},0)}+[3,0,0]_{(\frac{3}{2},0)}$ \\
\hline
$2+6$ & $[0,0,0]_{(0,0)+(0,2)+(1,1)+(2,0)}+[0,0,2]_{(\frac{1}{2},\frac{1}{2})+(\frac{1}{2},\frac{3}{2})}+[0,1,0]_{2(\frac{1}{2},\frac{1}{2})+(\frac{1}{2},\frac{3}{2})+(\frac{3}{2},\frac{1}{2})}+[2,0,2]_{(0,0)}+[2,1,0]_{(1,0)}$ \\
& $+[0,1,2]_{(0,1)}+[0,2,0]_{2(0,0)+(1,1)}+[1,0,1]_{(0,0)+2(0,1)+2(1,0)+(1,1)}+[1,1,1]_{2(\frac{1}{2},\frac{1}{2})}+[2,0,0]_{(\frac{1}{2},\frac{1}{2})+(\frac{3}{2},\frac{1}{2})}$ \\
\hline
$2+\frac{13}{2}$ & $[0,0,1]_{(0,\frac{1}{2})+(0,\frac{3}{2})+(1,\frac{1}{2})}+[0,1,1]_{(\frac{1}{2},0)+(\frac{1}{2},1)}+[1,0,0]_{(\frac{1}{2},0)+(\frac{1}{2},1)+(\frac{3}{2},0)}+[1,0,2]_{(0,\frac{1}{2})}$ \\
& $+[1,1,0]_{(0,\frac{1}{2})+(1,\frac{1}{2})}+[2,0,1]_{(\frac{1}{2},0)}$ \\
\hline
$2+7$ & $[0,0,0]_{(\frac{1}{2},\frac{1}{2})}+[0,0,2]_{(0,0)}+[0,1,0]_{(0,1)+(1,0)}+[1,0,1]_{(\frac{1}{2},\frac{1}{2})}+[2,0,0]_{(0,0)}$ \\
\hline
$2+\frac{15}{2}$ & $[0,0,1]_{(\frac{1}{2},0)}+[1,0,0]_{(0,\frac{1}{2})}$ \\
\hline
$2+8$ & $[0,0,0]_{(0,0)}$ \\
\hline
\end{longtable}
\normalsize

\end{document}
